% ================================================================
%  第 2 章  曲率估计及其推论
% ================================================================
\chapter{曲率估计及其推论}
\label{chap:2}

本章中, 我们将给出第 1 章所讨论的 Bernstein 定理的各种推广. 我们首先对
$\R^n$ 中极小超曲面 $\Sigma$ 的第二基本形式模长平方的 Laplace 导出 Simons
不等式. 在随后几节中, 我们将讨论这一不等式的各种应用. 第一个应用是 Choi 与
Schoen 的一个定理, 它给出总曲率很小的极小曲面的曲率估计. 利用这个估计, 我们
给出极小图的 Heinz 曲率估计的一个简短证明. 接下来, 我们讨论三维流形中稳定
极小曲面的先验估计, 包括 Colding 与 Minicozzi 关于面积与总曲率的估计, 以及
Schoen 的曲率估计. 之后, 我们遵循 Schoen、Simon 与 Yau 的做法, 把 Simons
不等式与稳定性不等式结合起来, 对稳定极小超曲面的第二基本形式模长平方证明
更高的 $L^p$ 界. 这些更高的 $L^p$ 界随后与 Simons 不等式一起用来证明稳定极小
超曲面的曲率估计, 并给出第 1 章所证 Bernstein 定理的一个推广, 它属于
De Giorgi、Almgren 与 Simons. 我们引入``几乎稳定''的概念, 它在理解嵌入曲面时
起关键作用. 接下来, 我们回到多值极小图, 证明 \cite{CM3} 中的一个重要结果:
若多值图有足够多的层, 则其间距的增长是次线性的. 本章最后讨论欧氏空间中的
极小锥, 以及 Bombieri、De Giorgi 与 Giusti 对维数大于七时 Bernstein 定理的
反例.

% ================================================================
\section{Simons 不等式}
\label{sec:2.1}

本节中, 我们将对 $\R^n$ 中极小超曲面 $\Sigma$ 的第二基本形式模长平方的
Laplace 导出一个非常有用的微分不等式. 在本章后面几节中, 我们将讨论这一
不等式的各种应用. 这个不等式最初属于 J. Simons \cite{Sim}; 它最一般的形式
断言: 对极小超曲面 $\Sigma^{n-1}\subset M^n$,
\begin{equation}
\Delta_{\Sigma}|A|^{2}\geq -C(1+|A|^{2})^{2},
\label{eq:2.1}
\end{equation}
其中 $C$ 依赖于 $M$ 的曲率及其协变导数.

下面我们来推导 Simons 不等式原来的形式; 上面那种更一般的不等式见
\cite{ScSiYa}.

\begin{lemma}[Simons 不等式, \cite{Sim}]
\label{lem:2.1}
设 $\Sigma^{n-1}\subset\R^n$ 是极小超曲面, 则
\begin{equation}
\Delta_{\Sigma}|A|^{2}\geq -2|A|^{4}
 +2\Bigl(1+\frac{2}{n-1}\Bigr)\bigl|\nabla_{\Sigma}|A|\bigr|^{2}.
\label{eq:2.2}
\end{equation}
\end{lemma}

注意这也可以等价地写成
\begin{equation}
|A|\Delta_{\Sigma}|A|+|A|^{4}
 \geq\frac{2}{n-1}\bigl|\nabla_{\Sigma}|A|\bigr|^{2}.
\label{eq:2.3}
\end{equation}

在证明该引理之前, 我们简要回忆 Riemann 曲率 $R$ 的定义, 以及用第二基本形式
表示 $R$ 的 Gauss 方程. 即: 若 $E_i$ 是流形 $\Sigma$(度量为 $g$, 联络为
$\nabla^{\Sigma}$)的标准正交标架, 则曲率 $R$ 由
\begin{equation}
\begin{aligned}
R_{ijkl}&=g\bigl(R(E_i,E_j)E_k,E_\ell\bigr)\\
&=g\bigl(\nabla^{\Sigma}_{E_j}\nabla^{\Sigma}_{E_i}E_k
 -\nabla^{\Sigma}_{E_i}\nabla^{\Sigma}_{E_j}E_k
 +\nabla^{\Sigma}_{\nabla^{\Sigma}_{E_i}E_j-\nabla^{\Sigma}_{E_j}E_i}E_k,
 E_\ell\bigr).
\end{aligned}
\label{eq:2.4}
\end{equation}
给出. 当 $\Sigma$ 是 $\R^n$ 中带诱导度量与诱导联络的超曲面时, Gauss 方程
用 $\Sigma$ 的第二基本形式表示曲率 $R$:
\begin{equation}
\mathrm{R}_{ijk\ell}=a_{ik}a_{j\ell}-a_{jk}a_{i\ell},
\label{eq:2.5}
\end{equation}
其中 $a_{jk}=\langle A(E_j,E_k),E_n\rangle$, 而 $E_n$ 是单位法向量.

\begin{proof}
设 $E_i$($i=1,\dots,n$)是某个 $x\in\Sigma$ 的邻域内局部定义的标准正交标架,
使得 $E_n$ 与 $\Sigma$ 正交. 令 $a$ 是 $\Sigma$ 上的对称二阶张量, 由
\begin{equation}
a(X,Y)=\langle A(X,Y),E_n\rangle=\langle\nabla_XY,E_n\rangle
 =-\langle\nabla_XE_n,Y\rangle
\label{eq:2.6}
\end{equation}
给出, 并置 $a_{ij}=-\langle\nabla_{E_i}E_n,E_j\rangle$. 也就是说,
\begin{equation}
a=\sum_{i,j=1}^{n-1}a_{ij}\Theta_i\otimes\Theta_j,
\label{eq:2.7}
\end{equation}
其中 $\Theta_i$ 是与 $E_i$ 对偶的标准正交余标架. 令 $a_{\cdot\cdot,k}$ 与
$a_{ij,k}$ 由
\begin{equation}
a_{\cdot\cdot,k}=\sum_{i,j=1}^{n-1}a_{ij,k}\Theta_i\otimes\Theta_j
 =\nabla_{E_k}a
\label{eq:2.8}
\end{equation}
定义. 由于 $a$ 是对称二阶张量, $a_{\cdot\cdot,k}$ 也是对称二阶张量. 由于
$\R^n$ 的曲率 $R$ 为零, 我们得到 Codazzi 方程:
\begin{equation}
\begin{aligned}
a_{ij,k}&=\bigl(\nabla^{T}_{E_k}a\bigr)(E_i,E_j)\\
&=E_ka(E_i,E_j)-a(\nabla^{T}_{E_k}E_i,E_j)-a(E_i,\nabla^{T}_{E_k}E_j)\\
&=-E_k\langle\nabla_{E_i}E_n,E_j\rangle
 +\langle\nabla_{\nabla^{T}_{E_k}E_i}E_n,E_j\rangle
 +\langle\nabla_{E_i}E_n,\nabla_{E_k}E_j\rangle\\
&=-\langle\nabla_{E_k}\nabla_{E_i}E_n,E_j\rangle
 +\langle\nabla_{\nabla^{T}_{E_k}E_i}E_n,E_j\rangle\\
\label{eq:2.9}
&=-\langle\nabla_{E_i}\nabla_{E_k}E_n,E_j\rangle
 +\langle\nabla_{\nabla^{T}_{E_i}E_k}E_n,E_j\rangle=a_{kj,i}.
\end{aligned}
\end{equation}
注意上面第二个等号就是 Leibniz 法则, 第三个等号是 $A$ 的定义, 最后一个等号
用到了 $R_{iknj}$ 在 $\R^n$ 上为零. 因此, 令 $a_{\cdot\cdot,\cdot}$ 是由
\begin{equation}
a_{\cdot\cdot,\cdot}=\sum_{i,j,k=1}^{n-1}a_{ij,k}
 \Theta_i\otimes\Theta_j\otimes\Theta_k
\label{eq:2.10}
\end{equation}
给出的对称三阶张量.

接下来, 我们定义对称三阶张量 $a_{\cdot\cdot,\ell}$ 为
\begin{equation}
a_{\cdot\cdot,\ell}=\sum_{i,j,k=1}^{n-1}a_{ij,k\ell}
 \Theta_i\otimes\Theta_j\otimes\Theta_k
 =\nabla_{E_\ell}a_{\cdot\cdot,\cdot}.
\label{eq:2.11}
\end{equation}

利用 Leibniz 法则, 我们算出
\begin{equation}
\begin{aligned}
a_{ij,k\ell}&=\nabla^{T}_{E_\ell}a_{\cdot\cdot,\cdot}(E_i,E_j,E_k)\\
&=E_\ell a_{\cdot\cdot,\cdot}(E_i,E_j,E_k)
 -a_{\cdot\cdot,\cdot}(\nabla^{T}_{E_\ell}E_i,E_j,E_k)\\
&\quad -a_{\cdot\cdot,\cdot}(E_i,\nabla^{T}_{E_\ell}E_j,E_k)
 -a_{\cdot\cdot,\cdot}(E_i,E_j,\nabla^{T}_{E_\ell}E_k)\\
&=E_\ell a_{\cdot\cdot,k}(E_i,E_j)
 -a_{\cdot\cdot,k}(\nabla^{T}_{E_\ell}E_i,E_j)\\
\label{eq:2.12}
&\quad -a_{\cdot\cdot,k}(E_i,\nabla^{T}_{E_\ell}E_j)
 -\nabla^{T}_{\nabla^{T}_{E_\ell}E_k}a(E_i,E_j).
\end{aligned}
\end{equation}

再应用一次 Leibniz 法则, 便得到
\begin{equation}
\begin{aligned}
a_{ij,k\ell}
&=E_\ell E_ka(E_i,E_j)-E_\ell a(\nabla^{T}_{E_k}E_i,E_j)
 -E_\ell a(E_i,\nabla^{T}_{E_k}E_j)\\
&\quad -E_ka(\nabla^{T}_{E_\ell}E_i,E_j)
 +a(\nabla^{T}_{E_k}\nabla^{T}_{E_\ell}E_i,E_j)
 +a(\nabla^{T}_{E_\ell}E_i,\nabla^{T}_{E_k}E_j)\\
&\quad -E_ka(E_i,\nabla^{T}_{E_\ell}E_j)
 +a(\nabla^{T}_{E_k}E_i,\nabla^{T}_{E_\ell}E_j)
 +a(E_i,\nabla^{T}_{E_k}\nabla^{T}_{E_\ell}E_j)\\
&\quad -\nabla^{T}_{E_\ell}E_k\,a(E_i,E_j)
 +a(\nabla^{T}_{\nabla^{T}_{E_\ell}E_k}E_i,E_j)
 +a(E_i,\nabla^{T}_{\nabla^{T}_{E_\ell}E_k}E_j)\\
\label{eq:2.13}
&=a_{ij,\ell k}+\sum_{m=1}^{n-1}R_{\ell kim}a_{mj}
 +\sum_{m=1}^{n-1}R_{\ell kjm}a_{mi}.
\end{aligned}
\end{equation}

(要看出最后一个等号, 请注意大多数项都成显然的对称对出现 —— 1 与 10、2 与 4、
3 与 7、6 与 8 —— 而剩下的项 —— 5 与 11、9 与 12 —— 给出曲率项.) 注意这一
计算对任何对称二阶张量 $a$ 都成立.

利用 Gauss 方程 \eqref{eq:2.5}, 我们可以把 \eqref{eq:2.13} 改写为
\begin{equation}
\begin{aligned}
a_{ik,jk}&=a_{ik,kj}+\sum_{m=1}^{n-1}\mathrm{R}_{kjim}\,a_{mk}
 +\sum_{m=1}^{n-1}\mathrm{R}_{kjkm}a_{mi}\\
&=a_{ik,kj}+\sum_{m=1}^{n-1}(a_{ki}a_{jm}-a_{ji}a_{km})a_{mk}\\
&\quad +\sum_{m=1}^{n-1}(a_{kk}a_{jm}-a_{jk}a_{km})a_{mi}.
\end{aligned}
\label{eq:2.14}
\end{equation}

最后, 我们准备证明 Simons 方程. 利用 \eqref{eq:2.14}, 得到
\begin{equation}
\begin{aligned}
\frac12\Delta_{\Sigma}|A|^{2}
&=\sum_{i,j=1}^{n-1}a_{ij}\Delta_{\Sigma}a_{ij}
 +\sum_{i,j=1}^{n-1}\bigl|\nabla_{\Sigma}a_{ij}\bigr|^{2}\\
&=\sum_{i,j,k=1}^{n-1}a_{ij}a_{ij,kk}+\sum_{i,j,k=1}^{n-1}a_{ij,k}^{2}
 =\sum_{i,j,k=1}^{n-1}a_{ij}a_{ik,jk}+\sum_{i,j,k=1}^{n-1}a_{ij,k}^{2}\\
&=\sum_{i,j,k=1}^{n-1}a_{ij}a_{kk,ij}
 +\sum_{i,j,k,m=1}^{n-1}a_{ij}(a_{ki}a_{jm}-a_{ji}a_{km})a_{mk}\\
&\quad +\sum_{i,j,k,m=1}^{n-1}a_{ij}(a_{kk}a_{jm}-a_{jk}a_{km})a_{mi}
 +\sum_{i,j,k=1}^{n-1}a_{ij,k}^{2}\\
&=-\sum_{i,j,k,m=1}^{n-1}a_{ij}^{2}a_{km}^{2}
 +\sum_{i,j,k=1}^{n-1}a_{ij,k}^{2}.
\end{aligned}
\label{eq:2.15}
\end{equation}

注意最后一个等号用到了 $\sum_ka_{kk}=0$(即 $\Sigma$ 是极小的), 并用 $a_{ij}$ 的
对称性消去了两个求和. 因此得到
\begin{equation}
\Delta_{\Sigma}|A|^{2}=-2|A|^{4}
 +2\sum_{i,j,k=1}^{n-1}a_{ij,k}^{2},
\label{eq:2.16}
\end{equation}
这就是 Simons 方程.

下面来看如何由此得到 Simons 不等式. 首先注意, 由于 $a$ 对称, 我们可以选取
$E_i$($i=1,\dots,n-1$), 使得在 $x$ 处有
\begin{equation}
a_{ij}=\lambda_i\delta_{ij}.
\label{eq:2.17}
\end{equation}
由于 $\nabla|A|^{2}=2|A|\nabla|A|$, 我们有
\begin{equation}
\begin{aligned}
4|A|^{2}\bigl|\nabla|A|\bigr|^{2}
&=\bigl|\nabla|A|^{2}\bigr|^{2}
 =\sum_{k=1}^{n-1}\Bigl[\Bigl(\sum_{i,j=1}^{n-1}a_{ij}^{2}\Bigr)_{k}\Bigr]^{2}\\
&=4\sum_{k=1}^{n-1}\Bigl(\sum_{i=1}^{n-1}a_{ii,k}\lambda_i\Bigr)^{2}
 \leq4|A|^{2}\sum_{i,k=1}^{n-1}a_{ii,k}^{2},
\end{aligned}
\label{eq:2.18}
\end{equation}
其中第三个等号用到了 \eqref{eq:2.17}, 而不等号用到了: 向量 $(a_{11,k},\dots,
a_{n-1,n-1,k})$ 与 $(\lambda_1,\dots,\lambda_{n-1})$ 的内积不超过它们长度之积.
于是 \eqref{eq:2.18} 给出
\begin{equation}
\bigl|\nabla|A|\bigr|^{2}\leq\sum_{i,k=1}^{n-1}a_{ii,k}^{2}.
\label{eq:2.19}
\end{equation}

因此, 由极小性与 \eqref{eq:2.19},
\begin{equation}
\begin{aligned}
\bigl|\nabla|A|\bigr|^{2}
&\leq\sum_{i,k=1}^{n-1}a_{ii,k}^{2}
 =\sum_{i\neq k}a_{ii,k}^{2}+\sum_{i=1}^{n-1}a_{ii,i}^{2}\\
&=\sum_{i\neq k}a_{ii,k}^{2}
 +\sum_{i=1}^{n-1}\Bigl(\sum_{i\neq j}a_{jj,i}\Bigr)^{2}\\
&\leq\sum_{i\neq k}a_{ii,k}^{2}
 +(n-2)\sum_{i=1}^{n-1}\sum_{i\neq j}a_{jj,i}^{2}
 =(n-1)\sum_{i\neq k}a_{ii,k}^{2}\\
&=(n-1)\sum_{i\neq k}a_{ik,i}^{2}
 =\frac{n-1}{2}\Bigl(\sum_{i\neq k}a_{ik,i}^{2}
 +\sum_{i\neq k}a_{ki,i}^{2}\Bigr).
\end{aligned}
\label{eq:2.20}
\end{equation}

\eqref{eq:2.20} 中的第二个不等号用到了代数事实
\begin{equation}
\Bigl(\sum_{j=1}^{n-2}b_j\Bigr)^{2}\leq(n-2)\sum_{j=1}^{n-2}b_j^{2}.
\label{eq:2.21}
\end{equation}

由 \eqref{eq:2.19} 与 \eqref{eq:2.20} 得到
\begin{equation}
\begin{aligned}
\Bigl(1+\frac{2}{n-1}\Bigr)\bigl|\nabla|A|\bigr|^{2}
&\leq\sum_{i,k=1}^{n-1}a_{ii,k}^{2}+\sum_{i\neq k}a_{ik,i}^{2}
 +\sum_{i\neq k}a_{ki,i}^{2}\\
&\leq\sum_{i,j,k=1}^{n-1}a_{ij,k}^{2}.
\end{aligned}
\label{eq:2.22}
\end{equation}
把 \eqref{eq:2.22} 与 Simons 方程(即 \eqref{eq:2.16})结合, 即得 Simons 不等式.
\end{proof}

\mtpara{曲面的 Simons 方程}
下面是该引理的一个推论:

若 $n=3$, $\Sigma^2\subset\R^3$ 是极小曲面, 则
\begin{equation}
\Delta\log|A|^{2}
 =\frac{\Delta|A|^{2}-4\bigl|\nabla|A|\bigr|^{2}}{|A|^{2}}
 =-2|A|^{2}.
\label{eq:2.23}
\end{equation}
这是因为, 对曲面而言 Simons 方程实际上给出的是等式
\begin{equation}
\Delta|A|^{2}=-2|A|^{4}+4\bigl|\nabla|A|\bigr|^{2},
\label{eq:2.24}
\end{equation}
而不只是不等式. 这很容易看出: 在曲面的情形, 上面用到的 Cauchy--Schwarz
不等式由极小性都取到等号.

方程 \eqref{eq:2.23} 有一个很好的几何解释:
\begin{equation}
|A|\langle\cdot,\cdot\rangle
\label{eq:2.25}
\end{equation}
是 $\Sigma$ 上的平坦(可能带奇点)度量. 即: 若 $(\Sigma,g)$ 是 Riemann 面, $f$
是 $\Sigma$ 上的正函数, 则 $(\Sigma,f^{2}g)$ 的曲率由
\begin{equation}
\mathrm{K}_g=\Delta_g\log f+f^{2}\mathrm{K}_{f^{2}g}
\label{eq:2.26}
\end{equation}
给出. 在此式中取 $f=|A|^{\frac12}$(其中 $|A|>0$), 并利用二维的 Simons
方程, 由极小性得到
\begin{equation}
\begin{aligned}
-|A|^{2}&=2\mathrm{K}_{\Sigma}
 =\Delta_{\Sigma}\log|A|+2|A|\mathrm{K}_{|A|\Sigma}\\
&=-|A|^{2}+2|A|\mathrm{K}_{|A|\Sigma}.
\end{aligned}
\label{eq:2.27}
\end{equation}
这显然说明 $|A|\langle\cdot,\cdot\rangle$ 在 $|A|>0$ 处是 $\Sigma$ 上的平坦
度量.

共形改变度量后曲率的公式(即公式 \eqref{eq:2.26})用活动标架很容易证明. 即: 对
$x\in\Sigma$, 设 $E_1$ 与 $E_2$ 是 $(\Sigma,g)$ 在 $x$ 的邻域内的标准正交标架,
$\Theta_i$ 是对偶的标准正交余标架. 则 $\bar\Theta_i=f\Theta_i$ 是
$(\Sigma,f^{2}g)$ 的标准正交余标架. 进一步, 由 Cartan 方程可见
\begin{equation}
\bar\omega_{1,2}=\omega_{1,2}-E_2(\log f)\Theta_1+E_1(\log f)\Theta_2,
\label{eq:2.28}
\end{equation}
从而
\begin{equation}
\begin{aligned}
\bar\Omega_{1,2}&=d\bar\omega_{1,2}
 =d\omega_{1,2}+\bigl[E_1(E_1(\log f))+E_2(E_2(\log f))\bigr]
 \Theta_1\wedge\Theta_2\\
&\quad -E_2(\log f)d\Theta_1+E_1(\log f)d\Theta_2\\
&=\Omega_{1,2}+\frac{\Delta_g\log f}{f^{2}}\bar\Theta_1\wedge\bar\Theta_2,
\end{aligned}
\label{eq:2.29}
\end{equation}
命题得证.

% ================================================================
\section{极小曲面的小能量曲率估计}
\label{sec:2.2}

本节中, 我们将对总曲率很小的极小曲面得到先验的内部曲率界. 这将是我们基本的
曲率估计, 后面在证明其他更一般的曲率估计时都会用到它, 例如极小图、稳定极小
曲面、面积很小的极小曲面等等.

下面我们证明 Choi 与 Schoen \cite{CiSc} 的``总曲率很小''估计. 它是几何分析中
一类非常重要的估计的一个例子. 可以与 J. Sacks 与 K. Uhlenbeck \cite{SaUh}
(对调和映射)、M. T. Anderson \cite{An} 以及 B. White \cite{Wh1} 的类似估计
作比较.

\begin{theorem}[Choi--Schoen, \cite{CiSc}]
\label{thm:2.2}
存在 $\epsilon,\rho>0$(依赖于 $M$), 使得: 若 $r_0<\rho$, $\Sigma^2\subset M$
是紧极小曲面且 $\partial\Sigma\subset\partial B_{r_0}(x)$, $0<\delta\leq1$, 并且
\begin{equation}
\int_{B_{r_0}\cap\Sigma}|A|^{2}<\delta\epsilon,
\label{eq:2.30}
\end{equation}
则对一切 $0<\sigma\leq r_0$ 与 $y\in B_{r_0-\sigma}(x)$,
\begin{equation}
\sigma^{2}|A|^{2}(y)\leq\delta.
\label{eq:2.31}
\end{equation}
\end{theorem}

\begin{proof}
我们只对 $M^n=\R^n$ 的情形给出证明, 一般情形所需的改动留给读者; 参看第 7 章.

在 $B_{r_0}\cap\Sigma$ 上置 $F=(r_0-r)^{2}|A|^{2}$, 注意 $F$ 在 $\partial B_{r_0}$
上为零. 设 $x_0\in\Sigma$ 是 $F$ 取到最大值的点, 并注意只需证明 $F(x_0)<\delta$.
我们将证明: 当 $\epsilon>0$ 足够小时, $F(x_0)\geq\delta$ 会导致矛盾.

由
\begin{equation}
\sigma^{2}|A|^{2}(x_0)=\frac{\delta}{4}
\label{eq:2.32}
\end{equation}
定义 $\sigma>0$. 特别地, 由于 $F(x_0)\geq\delta$, 可见
\begin{equation}
2\sigma\leq r_0-r(x_0).
\label{eq:2.33}
\end{equation}
于是由三角不等式, 在 $B_\sigma(x_0)$ 上有
\[
\frac12\leq\frac{r_0-r}{r_0-r(x_0)}\leq2.
\]
由于 $F$ 在 $x_0$ 取到最大值, 我们断定
\begin{equation}
\begin{aligned}
(r_0-r(x_0))^{2}\sup_{B_\sigma(x_0)\cap\Sigma}|A|^{2}
&\leq4\sup_{B_\sigma(x_0)\cap\Sigma}F(x)=4F(x_0)\\
&=4(r_0-r(x_0))^{2}|A|^{2}(x_0).
\end{aligned}
\label{eq:2.34}
\end{equation}
两边除以 $(r_0-r(x_0))^{2}$ 并利用 $\sigma$ 的定义, 得到
\begin{equation}
\sup_{B_\sigma(x_0)\cap\Sigma}|A|^{2}\leq4|A|^{2}(x_0)=\delta\sigma^{-2}.
\label{eq:2.35}
\end{equation}

把 $B_\sigma(x_0)\cap\Sigma$ 伸缩到单位尺寸(仍记作 $\Sigma$!)后, 我们有
\begin{equation}
\sup_{B_1(x_0)\cap\Sigma}|A|^{2}\leq4|A|^{2}(x_0)=\delta\leq1.
\label{eq:2.36}
\end{equation}
由 $B_1(x_0)\cap\Sigma$ 上的 Simons 不等式,
\begin{equation}
\Delta|A|^{2}\geq-2|A|^{2}.
\label{eq:2.37}
\end{equation}
所要的矛盾现在由平均值不等式得出. 即: 推论~\ref{cor:1.16} 表明
\begin{equation}
\frac{\delta}{4}=|A|^{2}(x_0)
 \leq\mathrm{e}\,\frac{\int_{B_1(x_0)}|A|^{2}}{\pi}
 <\frac{\mathrm{e}}{\pi}\delta\epsilon,
\label{eq:2.38}
\end{equation}
只要 $\epsilon$ 取得足够小, 这便是矛盾.
\end{proof}

\mtpara{极小图的 Heinz 曲率估计}
作为一个直接应用, 下面我们证明 E. Heinz 关于极小图的曲率估计. 为此, 设
$D_{r_0}$ 是 $\R^2$ 中半径为 $r_0$ 的圆盘.

\begin{theorem}[Heinz, \cite{He}]
\label{thm:2.3}
若 $u:D_{r_0}\to\R$ 满足极小曲面方程, 则对 $\Sigma=\Graph_u$ 与
$0<\sigma\leq r_0$,
\begin{equation}
\sigma^{2}\sup_{D_{r_0-\sigma}}|A|^{2}\leq C.
\label{eq:2.39}
\end{equation}
\end{theorem}

\begin{proof}
显然只需证明 $\sigma=r_0$ 的情形. 在上一章的推论~\ref{cor:1.20} 中我们已经
证明: 对任意 $\kappa>1$,
\begin{equation}
\int_{B_{r_0/\sqrt{\kappa}}\cap\Graph_u}|A|^{2}
 \leq\frac{C}{\log\kappa}.
\label{eq:2.40}
\end{equation}
我们选取 $\kappa>1$ 使得
\begin{equation}
\frac{C}{\log\kappa}<\epsilon,
\label{eq:2.41}
\end{equation}
其中 $\epsilon$ 由定理~\ref{thm:2.2} 给出, 于是定理立即得证.
\end{proof}

到此我们已经看到, 总曲率很小的极小曲面与极小图都满足先验曲率估计. 下面我们
还会看到, 面积很小也蕴含曲率估计(见定理~\ref{thm:2.15}).

\mtpara{内蕴球的估计}
为此, 方便的做法是用 $B_r^{\Sigma}(x)$ 表示 $\Sigma$ 中以 $x$ 为心、半径为 $r$
的内蕴球:
\begin{equation}
B_r^{\Sigma}(x)=\bigl\{y\in\Sigma\mid\dist_{\Sigma}(x,y)<r\bigr\}.
\label{eq:2.42}
\end{equation}
显然总有 $B_r^{\Sigma}(x)\subset B_r(x)$. 本节余下部分都用这个记号.

我们将用到下面这个初等引理, 它说明: 在有曲率界时, $\Sigma$ 中小的内蕴球是
图形化的.

\begin{lemma}[小曲率蕴含图形化]
\label{lem:2.4}
设 $\Sigma^2\subset\R^n$ 是浸入曲面, 且
\begin{equation}
16s^{2}\sup_{\Sigma}|A_{\Sigma}|^{2}\leq1.
\label{eq:2.43}
\end{equation}
若 $x\in\Sigma$ 且 $\dist_{\Sigma}(x,\partial\Sigma)\geq2s$, 则

\begin{enumerate}
\item $B_{2s}^{\Sigma}(x)$ 可以写成 $T_x\Sigma$ 上某个函数 $u$ 的图, 且
  $|\nabla u|\leq1$, $\sqrt{2}s|\mathrm{Hess}_u|\leq1$;
\item $B_s(x)\cap\Sigma$ 中含 $x$ 的连通分支包含于 $B_{2s}^{\Sigma}(x)$.
\end{enumerate}
\end{lemma}

引理的第 (2) 部分是一个弦弧界, 它把内蕴距离与外蕴距离联系起来.

\begin{proof}
在本引理的证明中, 我们定义
\begin{equation}
d_{x,y}\equiv\dist_{\bS^{n-1}}\bigl(N(x),N(y)\bigr).
\label{eq:2.44}
\end{equation}
回忆由 \eqref{eq:1.84} 有 $|\nabla N|\leq|A|$. 因此, 给定 $y\in\Sigma$, 沿连接 $x$ 与
$y$ 的测地线对 \eqref{eq:2.43} 积分, 得到
\begin{equation}
\sup_{y\in B_{2s}^{\Sigma}(x)}d_{x,y}\leq\frac12<\frac{\pi}{4}.
\label{eq:2.45}
\end{equation}
由于在此集合上 $d_{x,y}<\pi/2$, 可见 $B_{2s}^{\Sigma}(x)$ 含于 $T_x\Sigma$
的(某个子集)上某个函数 $u$ 的图中. 此外, 利用 \eqref{eq:1.3}, 我们有
\begin{equation}
1+|\nabla u|^{2}=\langle N(x),N(y)\rangle^{-2}=\cos^{-2}d_{x,y},
\label{eq:2.46}
\end{equation}
于是 \eqref{eq:2.45} 蕴含 $|\nabla u|\leq1$. 对 $u$ 的 Hessian 估计由梯度界、曲率界以及
\eqref{eq:1.94} 给出:
\begin{equation}
|\mathrm{Hess}_u|^{2}\leq(1+|\nabla u|^{2})^{2}|A|^{3}
 \leq8\frac{1}{16s^{2}}.
\label{eq:2.47}
\end{equation}

最后, 为证明 $B_s(x)\cap\Sigma$ 中含 $x$ 的连通分支包含于 $B_{2s}^{\Sigma}(x)$,
我们将证明 $\partial B_{2s}^{\Sigma}(x)$ 落在 $\overline{B_s(x)}$ 之外. 由
初等几何, 我们知道 $\partial B_{2s}^{\Sigma}(x)$ 上每一点都是 $\Sigma$ 中一条
从 $x$ 出发、长度为 $2s$ 的内蕴测地线的端点. 设
\[
\gamma:[0,2s]\to\Sigma
\]
是这样一条按弧长参数化的测地线. 注意
\begin{equation}
\bigl|\partial_t\langle\gamma'(t),\gamma'(0)\rangle\bigr|
 \leq|A|(\gamma(t))\leq1/4s.
\label{eq:2.48}
\end{equation}
由于 $\langle\gamma'(0),\gamma'(0)\rangle=1$, 对任意 $t\leq2s$ 把它从 $0$ 到
$t$ 积分, 得到
\begin{equation}
\langle\gamma'(t),\gamma'(0)\rangle\geq1-\frac{t}{4s}\geq\frac12.
\label{eq:2.49}
\end{equation}
再把它从 $0$ 到 $2s$ 积分, 得到
\begin{equation}
\bigl\langle[\gamma(t)-\gamma(0)],\gamma'(0)\bigr\rangle
 \geq\int_0^{2s}\Bigl(1-\frac{t}{4s}\Bigr)\,dt=\frac{3s}{2}.
\label{eq:2.50}
\end{equation}
由此可见 $\gamma(2s)$ 落在 $B_{3s/2}(x)$ 之外. 由于这对任何这样的 $\gamma$ 都
成立, 我们得到
\begin{equation}
B_{3s/2}(x)\cap\partial B_{2s}^{\Sigma}(x)=\emptyset,
\label{eq:2.51}
\end{equation}
证明完成.
\end{proof}

注意, 若把 \eqref{eq:2.43} 右端的界 $1$ 换成某个 $\delta<1$, 则我们会得到
\begin{equation}
\sup_{y\in\tilde\Sigma}d_{x,y}\leq\frac{\delta}{2}.
\label{eq:2.52}
\end{equation}
若 $0<s<\frac{\pi}{4}$, 则
\begin{equation}
\frac{\partial}{\partial s}\cos^{-2}s=2\sin s\cos^{-3}s\leq4,
\label{eq:2.53}
\end{equation}
于是 $\cos^{-2}s\leq1+4s$. 因此, \eqref{eq:2.46} 与 \eqref{eq:2.52} 蕴含
\begin{equation}
|\nabla u|^{2}\leq4\delta.
\label{eq:2.54}
\end{equation}

现在可以看出:

\begin{remark}
\label{rem:2.5}
定理~\ref{thm:2.2} 用内蕴球代替外蕴球仍然成立(常数略有不同).
\end{remark}

要看出这一点, 改用内蕴距离记作 $r$, 把外蕴球都换成内蕴球, 并逐字重复论证直到
\eqref{eq:2.37}. 此时, 按引理~\ref{lem:2.4} 中的论证, \eqref{eq:2.36} 蕴含 $B_1^{\Sigma}(x_0)$
包含 $B_{1/2}(x_0)$ 中含 $x_0$ 的连通分支. 现在我们可以把 $B_1(x_0)$ 换成
$B_{1/2}(x_0)$, 并像前面一样应用平均值不等式来完成证明.

\mtpara{高阶导数的估计}
由引理~\ref{lem:2.4}, 前一结果的曲率界 \eqref{eq:2.31} 蕴含 $B_{\sigma/4}(y)\cap\Sigma$
的连通分支是图形化的, 且梯度与 Hessian 都有界. 这说明相应的极小曲面方程是
一个具有 $C^{1}$ 系数的、一致椭圆的散度型方程. 标准的椭圆理论于是给出统一的
估计(例如可参看 \cite{GiTr} 的推论 16.7). 第 7 章研究极小曲面的紧性定理时
将会用到这一点.

% ================================================================
\section{曲率与面积}
\label{sec:2.3}

本节中, 我们将证明若干依赖于面积与总曲率估计的二维结果. 基本想法很简单, 但
仍然相当有用. 本节的结果取自 \cite{CM2} 与 \cite{CM4}.

出发点是下面这个联系面积增长与总曲率的引理:

\begin{lemma}
\label{lem:2.6}
若 $B_{r_0}^{\Sigma}(x)\subset\Sigma^2$ 与 $x$ 的割迹不相交, 则
\begin{equation}
\Length(\partial B_{r_0}^{\Sigma})-2\pi r_0
 =-\int_0^{r_0}\int_{B_\rho^{\Sigma}}\mathrm{K}_{\Sigma},
\label{eq:2.55}
\end{equation}
\begin{equation}
\Area\bigl(B_{r_0}^{\Sigma}(x)\bigr)-\pi r_0^{2}
 =-\int_0^{r_0}\int_0^{\tau}\int_{B_\rho^{\Sigma}(x)}\mathrm{K}_{\Sigma}.
\label{eq:2.56}
\end{equation}
\end{lemma}

\begin{proof}
对 $0<t\leq r_0$, 由 Gauss--Bonnet 定理(参看 \cite{dC1} 第 274 页),
\begin{equation}
\frac{d}{dt}\int_{\partial B_t^{\Sigma}}1
 =\int_{\partial B_t^{\Sigma}}k_g
 =2\pi-\int_{B_t^{\Sigma}}\mathrm{K}_{\Sigma},
\label{eq:2.57}
\end{equation}
其中 $k_g$ 是 $\partial B_t^{\Sigma}$ 的测地曲率. 对 \eqref{eq:2.57} 积分即得引理.
\end{proof}

下面这个取自 \cite{CM4} 的推论把引理~\ref{lem:2.6} 特殊化到 $\Sigma$ 极小的
情形.

\begin{corollary}
\label{cor:2.7}
若 $\Sigma^2\subset\R^3$ 是浸入极小曲面, $B_{r_0}^{\Sigma}\subset\Sigma^2$ 是
圆盘, 且 $B_{r_0}^{\Sigma}\cap\partial\Sigma=\emptyset$, 则
\begin{equation}
\begin{aligned}
t^{2}\int_{B_{r_0-2t}^{\Sigma}}|A|^{2}
&\leq\int_{B_{r_0}^{\Sigma}}|A|^{2}(r_0-r)^{2}/2
 =\int_0^{r_0}\int_0^{\tau}\int_{B_\rho^{\Sigma}(x)}|A|^{2}\\
&=2\bigl(\Area(B_{r_0}^{\Sigma})-\pi r_0^{2}\bigr)
 \leq r_0\Length(\partial B_{r_0}^{\Sigma})-2\pi r_0^{2}.
\end{aligned}
\label{eq:2.58}
\end{equation}
\end{corollary}

\begin{proof}
由于 $\Sigma$ 极小, Gauss 方程给出 $|A|^{2}=-2\mathrm{K}_{\Sigma}$, 于是由
引理~\ref{lem:2.6} 得到
\begin{equation}
\begin{aligned}
t^{2}\int_{B_{r_0-2t}^{\Sigma}}|A|^{2}
&\leq t\int_0^{r_0-t}\int_{B_\rho^{\Sigma}}|A|^{2}
 \leq\int_0^{r_0}\int_0^{\tau}\int_{B_\rho^{\Sigma}}|A|^{2}\\
&=2\bigl(\Area(B_{r_0}^{\Sigma})-\pi r_0^{2}\bigr).
\end{aligned}
\label{eq:2.59}
\end{equation}
\eqref{eq:2.58} 中的第一个等号由余面积公式与两次分部积分得到:
\begin{equation}
\int_0^{r_0}f(t)g''(t)\,dt=\int_0^{r_0}f''(t)g(t)\,dt,
\label{eq:2.60}
\end{equation}
其中 $f(t)=\int_0^t\int_{B_s^{\Sigma}}|A|^{2}$, $g(t)=(r_0-t)^{2}/2$.

为得到 \eqref{eq:2.58} 中最后一个不等号, 注意由 \eqref{eq:2.57} 有
$\frac{d^{2}}{dt^{2}}\Length(\partial B_t^{\Sigma})\geq0$, 从而
\[
t\frac{d}{dt}\Length(\partial B_t^{\Sigma})\geq\Length(\partial B_t^{\Sigma}),
\]
于是
\[
\frac{d}{dt}\Bigl(\Length(\partial B_t^{\Sigma})/t\Bigr)\geq0.
\]
由此容易得到所需结论.
\end{proof}

\mtpara{稳定极小曲面的估计}
对三维流形中的稳定极小曲面, 要得到曲率估计并不需要关于总曲率、面积或拓扑的
任何假设. 在高维情形这一点尚不知道. 为简单与清晰起见, 我们把讨论限制在
$\R^3$ 中的曲面.

出发点是 $\R^3$ 中双边稳定极小曲面的一个统一面积界:

\begin{theorem}[Colding--Minicozzi, \cite{CM2}]
\label{thm:2.8}
若 $\Sigma^2\subset\R^3$ 稳定、双边, 且 $B_{r_0}^{\Sigma}$ 单连通, 则
\begin{equation}
\Area\bigl(B_{r_0}^{\Sigma}\bigr)\leq\frac{4\pi}{3}r_0^{2}.
\label{eq:2.61}
\end{equation}
\end{theorem}

\begin{proof}
设 $r$ 是到 $B_{r_0}^{\Sigma}$ 中心的内蕴距离函数. 先用推论~\ref{cor:2.7}, 再用
稳定性不等式(引理~\ref{lem:1.32}), 得到
\begin{equation}
\begin{aligned}
4\bigl(\Area(B_{r_0}^{\Sigma})-\pi r_0^{2}\bigr)
&=\int_{B_{r_0}^{\Sigma}}|A|^{2}(r_0-r)^{2}
 \leq\int_{B_{r_0}^{\Sigma}}|\nabla r|^{2}\\
&=\Area\bigl(B_{r_0}^{\Sigma}\bigr),
\end{aligned}
\label{eq:2.62}
\end{equation}
其中最后一个等号用到了 $|\nabla r|=1$. 定理得证.
\end{proof}

作为一个直接推论, 我们可以在定理~\ref{thm:2.8} 中令 $r_0\to\infty$, 从而得到
Bernstein 型定理(参看 \cite{FiSc} 或 \cite{dCP}):

\begin{corollary}
\label{cor:2.9}
若 $\Sigma^2\subset\R^3$ 完全、稳定且双边, 则它是平坦的.
\end{corollary}

\begin{proof}
由命题~\ref{prop:1.39}, $\Sigma$ 的万有覆盖也稳定, 故可设 $\Sigma$ 单连通.
由于 $\Sigma$ 具有非正曲率(由 Gauss 方程), $\Sigma$ 中的测地球单连通, 于是
定理~\ref{thm:2.8} 蕴含 $\Sigma$ 具有二次面积增长. 命题~\ref{prop:1.37} 随后
蕴含 $\Sigma$ 是抛物型的. 最后, 推论~\ref{cor:1.40} 蕴含 $\Sigma$ 是平坦的.
\end{proof}

仿照 \cite{CM2}, 我们可以用定理~\ref{thm:2.8} 证明 R. Schoen 的下述估计:

\begin{theorem}[\cite{Sc1}]
\label{thm:2.10}
若 $\Sigma^2\subset M^3$ 是浸入稳定极小曲面, 法丛平凡,
$B_{r_0}^{\Sigma}=B_{r_0}^{\Sigma}(x)\subset\Sigma\setminus\partial\Sigma$,
其中 $|K_M|\leq k^{2}$ 且 $r_0<\rho_1(\pi/k,k)$, 则存在 $C=C(k)$, 使得对一切
$0<\sigma\leq r_0$,
\begin{equation}
\sup_{B_{r_0-\sigma}^{\Sigma}}|A|^{2}\leq C\sigma^{-2}.
\label{eq:2.63}
\end{equation}
\end{theorem}

\begin{proof}
我们对 $M=\R^3$ 给出 \cite{CM2} 中的证明, 容易的改动留给读者. 我们先作两点
观察. 第一, 只需对 $\sigma=r_0$ 证明 \eqref{eq:2.63}, 即证明
\begin{equation}
|A|^{2}(x)\leq Cr_0^{-2}.
\label{eq:2.64}
\end{equation}
第二, 不妨设 $\Sigma$ 是拓扑圆盘, 因为命题~\ref{prop:1.39} 蕴含 $\Sigma$ 的
万有覆盖也稳定.

由于 $\Sigma$ 具有非正曲率且可设它是圆盘, 定理~\ref{thm:2.8} 给出: 对每个
$s\leq r_0$, $\Area(B_s^{\Sigma}(x))\leq\frac{4\pi}{3}s^{2}$. 这个二次面积界
使我们能用对数截断技巧在某个子球上得到很小的总曲率. 即: 设 $r$ 表示 $\Sigma$
中到 $x$ 的距离, 固定一个待定的大整数 $n>0$, 并定义 $\eta$ 为
\begin{equation}
\eta=\begin{cases}
1 & \text{若 } r\leq\mathrm{e}^{-n}r_0,\\
\dfrac{\log r_0-\log r}{n}
 & \text{若 } \mathrm{e}^{-n}r_0<r\leq r_0,\\
0 & \text{若 } r>r_0.
\end{cases}
\label{eq:2.65}
\end{equation}
由于 $|\nabla r|=1$, 我们有 $|\nabla\eta|\leq\frac{1}{nr}$, 且 $|\nabla\eta|$ 的
支集位于 $\mathrm{e}^{-n}r_0$ 与 $r_0$ 之间的环域内. 把稳定性不等式(即
引理~\ref{lem:1.32})用于截断函数 $\eta$, 并利用上面的面积界, 得到
\begin{equation}
\begin{aligned}
\int_{B_{\mathrm{e}^{-n}r_0}^{\Sigma}\cap\Sigma}|A|^{2}
&\leq\int_\Sigma\eta^{2}|A|^{2}\leq\int_\Sigma|\nabla\eta|^{2}
 \leq n^{-2}\int_{B_{r_0}^{\Sigma}\setminus
 B_{\mathrm{e}^{-n}r_0}^{\Sigma}}r^{-2}\\
&\leq n^{-2}\sum_{\ell=-n}^{-1}
 \int_{\bigl(B_{\mathrm{e}^{\ell+1}r_0}^{\Sigma}
 \setminus B_{\mathrm{e}^{\ell}r_0}^{\Sigma}\bigr)}r^{-2}\\
&\leq n^{-2}\sum_{\ell=-n}^{-1}\frac{4}{3}\pi\mathrm{e}^{2}
 =\frac{4\pi\mathrm{e}^{2}}{3n}.
\end{aligned}
\label{eq:2.66}
\end{equation}
因此, 当 $n$ 足够大时, 内蕴版本的 Choi--Schoen 小总曲率估计(定理~\ref{thm:2.2}
与注~\ref{rem:2.5})给出 \eqref{eq:2.64}.
\end{proof}

由于内蕴球包含于外蕴球:

\begin{corollary}[R. Schoen, \cite{Sc1}]
\label{cor:2.11}
若 $\Sigma^2\subset B_{r_0}=B_{r_0}(x)\subset M^3$ 是浸入稳定极小曲面, 法丛
平凡, 其中 $|K_M|\leq k^{2}$, $r_0<\rho_1(\pi/k,k)$, 且
$\partial\Sigma\subset\partial B_{r_0}$, 则存在 $C=C(k)$, 使得对一切
$0<\sigma\leq r_0$,
\begin{equation}
\sup_{B_{r_0-\sigma}}|A|^{2}\leq C\sigma^{-2}.
\label{eq:2.67}
\end{equation}
\end{corollary}

R. Schoen 猜测这对 $\R^4$ 中的稳定超曲面也应当成立:

\begin{conjecture}[Schoen]
\label{conj:2.12}
若 $\Sigma^3\subset\R^4$ 是完全浸入稳定极小超曲面, 法丛平凡, 则 $\Sigma$ 是
平坦的.
\end{conjecture}

\begin{conjecture}[Schoen]
\label{conj:2.13}
若 $\Sigma^3\subset B_{r_0}=B_{r_0}(x)\subset M^4$ 是浸入稳定极小超曲面, 法丛
平凡, 其中 $|K_M|\leq k^{2}$, $r_0<\rho_1(\pi/k,k)$, 且
$\partial\Sigma\subset\partial B_{r_0}$, 则存在 $C=C(k)$, 使得对一切
$0<\sigma\leq r_0$,
\begin{equation}
\sup_{B_{r_0-\sigma}}|A|^{2}\leq C\sigma^{-2}.
\label{eq:2.68}
\end{equation}
\end{conjecture}

把定理~\ref{thm:2.10} 推广到对某个参数型椭圆被积函数稳定的曲面, 见
\cite{CM2}. 在继续之前, 我们简要讨论这一推广.

设 $M^3$ 是 Riemann 三维流形. 给定 $M$ 的单位球丛上的函数 $\phi\geq1$, 我们
可以在 $M$ 中浸入定向曲面 $\Sigma$ 上定义泛函 $\Phi$ 为
\begin{equation}
\Phi(\Sigma)=\int_{x\in\Sigma}\phi\bigl(x,N(x)\bigr)\,dx.
\label{eq:2.69}
\end{equation}
限制 $\phi\geq1$ 是一个方便的规范化. 与面积泛函(即 $\phi\equiv1$ 的情形)类比,
若曲面是 $\Phi$ 的临界点, 则称它是 \textbf{$\Phi$-驻定}的; 若 $\Phi$-驻定
曲面的第二变分对紧支集变分非负, 则称它是 \textbf{$\Phi$-稳定}的. 关于参数型
被积函数的第一、第二变分, 进一步的讨论见 \cite{Al2}. 若存在 $\lambda>0$, 使得
对每个 $x\in M$, $v\mapsto[\phi(x,v/|v|)-\lambda]|v|$ 是 $v\in T_xM$ 的凸函数,
则称 $\Phi$ 是\textbf{椭圆}的.

对 $\Phi$-稳定曲面, 曲率估计仍然成立. 即: 设 $D_1\phi$ 与 $D_2\phi$ 分别表示
$\phi$ 对第一个与第二个分量的导数; 则有

\begin{theorem}[Colding--Minicozzi, \cite{CM2}]
\label{thm:2.14}
存在 $\epsilon,\rho>0$ 与 $C$, 其中 $\rho,C$ 依赖于
$|\phi|_{C^{2,\alpha}}$ 与 $|D_2\phi|_{C^{2,\alpha}}$, 使得: 若 $r_0\leq\rho$,
$|D_2\phi|+|D_2^{2}\phi|<\epsilon$, 且 $B_{r_0}^{\Sigma}=B_{r_0}^{\Sigma}(x)$
是浸入定向 $\Phi$-稳定曲面 $\Sigma\subset\R^3$ 中的球, 满足
$B_{r_0}^{\Sigma}\cap\partial\Sigma=\emptyset$, 则对一切 $0<\sigma\leq r_0$,
\begin{equation}
\sup_{B_{r_0-\sigma}^{\Sigma}}|A|^{2}\leq C\sigma^{-2}.
\label{eq:2.70}
\end{equation}
\end{theorem}

与极小曲面方程的情形一样, 标准椭圆理论随后给出高阶导数的估计(参看
\cite{Si2}; 并参看 \cite{GiTr} 第 16 章).

我们指出, 极小曲面的许多标准工具对一般的 $\Phi$ 并不成立. 特别地, 面积的
单调性(即命题~\ref{prop:1.12})与 Simons 不等式(即引理~\ref{lem:2.1})都不再
成立. 这带来了相当大的麻烦. 例如, \cite{Sc1} 中就用 Simons 不等式与一个迭代
论证来估计稳定浸入的共形因子.

\mtpara{小超出量曲率估计}
下面我们仿照 \cite{CM2}, 对具有小\newterm{excess}{超出量}{excess}的(光滑)
极小曲面证明曲率估计. Allard 的正则性定理 \cite{Al1} 表明, 这一估计在更一般的
情形下也成立, 而不必假设光滑性.

设 $\Sigma^2\subset B_1$ 是光滑极小曲面, $\partial\Sigma\subset\partial B_1$.
给定 $0<s<1$ 与 $x\in B_{1-s}\cap\Sigma$, $\Sigma$ 在 $B_s(x)$ 中的超出量定义
为
\begin{equation}
\Theta_x(s)-1=\frac{\Area(B_s(x)\cap\Sigma)}{\pi s^{2}}-1.
\label{eq:2.71}
\end{equation}
单调性(即命题~\ref{prop:1.12})蕴含超出量非负, 且关于 $s$ 单调递增.

\begin{theorem}
\label{thm:2.15}
存在 $\epsilon>0$ 与 $C<\infty$, 使得: 若
$0\in\Sigma^2\subset B_{2r_0}$ 是光滑紧极小曲面,
$\partial\Sigma\subset\partial B_{2r_0}$, 且对一切 $y\in B_{r_0}\cap\Sigma$,
\begin{equation}
\Theta_y(r_0)-1<\epsilon,
\label{eq:2.72}
\end{equation}
则对一切 $0<\sigma\leq r_0$ 与 $y\in B_{r_0-\sigma}(x)$,
\begin{equation}
\sigma^{2}|A|^{2}(y)\leq1.
\label{eq:2.73}
\end{equation}
\end{theorem}

\begin{proof}
首先注意, 由单调性, \eqref{eq:2.72} 蕴含: 对一切 $s<r_0$ 与 $y\in B_{r_0}\cap\Sigma$,
\begin{equation}
\Theta_y(s)-1<\epsilon.
\label{eq:2.74}
\end{equation}
在 $B_{r_0}\cap\Sigma$ 上定义 $F=(r_0-r)^{2}|A|^{2}$, 只需证明 $F\leq1$. 若
不然, 按定理~\ref{thm:2.2} 的证明中那样论证, 于是在伸缩 $\Sigma$ 之后有
\begin{equation}
|A|^{2}(x_0)=\frac14
\label{eq:2.75}
\end{equation}
以及
\begin{equation}
\sup_{B_1(x_0)\cap\Sigma}|A|^{2}\leq1.
\label{eq:2.76}
\end{equation}
此外, \eqref{eq:2.74} 蕴含对任意 $s<1$,
\begin{equation}
\Area(B_s(x_0)\cap\Sigma)\leq\pi(1+\epsilon)s^{2}.
\label{eq:2.77}
\end{equation}
为完成证明, 我们将证明当 $\epsilon>0$ 足够小时, \eqref{eq:2.75}、\eqref{eq:2.76} 与 \eqref{eq:2.77} 会
导致矛盾.

首先, 应用定理~\ref{thm:2.2}, \eqref{eq:2.75} 蕴含
\begin{equation}
\pi C\leq\int_{B_{\frac{1}{32}}(x_0)\cap\Sigma}|A|^{2},
\label{eq:2.78}
\end{equation}
其中 $C$ 来自定理~\ref{thm:2.2}. 由引理~\ref{lem:2.4}, $B_{1/8}(x_0)\cap\Sigma$
是梯度不超过 $1$ 的图. 特别地, 它单连通, 且对 $x,y\in B_{1/8}(x_0)\cap\Sigma$,
\begin{equation}
|x-y|\leq\dist_{\Sigma}(x,y)\leq\sqrt{2}|x-y|.
\label{eq:2.79}
\end{equation}
在本证明余下部分, $B_s^{\Sigma}$ 表示 $\Sigma$ 中以 $x_0$ 为心、半径为 $s$ 的
内蕴球. 结合 \eqref{eq:2.78} 与 \eqref{eq:2.79}, 得到
\begin{equation}
\pi C\leq\int_{B_{1/16}^{\Sigma}}|A|^{2}.
\label{eq:2.80}
\end{equation}
把推论~\ref{cor:2.7} 用于 $r_0=1/8$ 与 $t=1/32$, 得到
\begin{equation}
(1/32)^{2}\int_{B_{1/16}^{\Sigma}}|A|^{2}
 \leq2\bigl(\Area(B_{1/8}^{\Sigma})-\pi(1/8)^{2}\bigr),
\label{eq:2.81}
\end{equation}
于是
\begin{equation}
\Area(B_{\frac18}^{\Sigma})\geq\pi\Bigl(1+\frac{C}{32}\Bigr)\frac{1}{64}.
\label{eq:2.82}
\end{equation}
最后, \eqref{eq:2.79} 与 \eqref{eq:2.77} 给出
\begin{equation}
\Area(B_{\frac18}^{\Sigma})
 \leq\Area(B_{\frac18}(x_0)\cap\Sigma)\leq\frac{\pi(1+\epsilon)}{64},
\label{eq:2.83}
\end{equation}
当 $\epsilon$ 足够小时这与 \eqref{eq:2.82} 矛盾. 此矛盾完成证明.
\end{proof}

上面的证明表明, 对曲面而言, 定理~\ref{thm:2.2} 蕴含 Allard 估计的光滑版本.
事实上可以证明, 用类似的论证可知这两个估计是等价的. 关于这一点及相关结果,
\cite{CM2} 中有更详细的讨论.

\mtpara{单连通曲面的曲率估计}
事实上, 对单连通的嵌入极小曲面, 关于面积或总曲率的任何界都蕴含曲率估计. 这
对具有拟共形 Gauss 映射的曲面(参看 \eqref{eq:7.6})更一般地成立, 但我们只限于极小曲面.

这类结果中的第一个是 Schoen 与 Simon 关于嵌入极小圆盘中外蕴球的估计:

\begin{theorem}[Schoen--Simon, \cite{ScSi}]
\label{thm:2.16}
设 $0\in\Sigma^2\subset B_{r_0}=B_{r_0}(0)\subset\R^3$ 是嵌入单连通极小曲面,
$\partial\Sigma\subset\partial B_{r_0}$. 若 $\mu>0$, 且或者
\begin{equation}
\Area(\Sigma)\leq\mu r_0^{2},\quad\text{或}
\label{eq:2.84}
\end{equation}
\begin{equation}
\int_\Sigma|A|^{2}\leq\mu,
\label{eq:2.85}
\end{equation}
则对 $B_{r_0/2}\cap\Sigma$ 中含 $0$ 的连通分支 $\Sigma'$, 有
\begin{equation}
\sup_{\Sigma'}|A|^{2}\leq Cr_0^{-2}
\label{eq:2.86}
\end{equation}
其中 $C=C(\mu)$.
\end{theorem}

作为一个直接推论, 令 $r_0\to\infty$ 便得到密度有界或总曲率有限的嵌入单连通
极小曲面的 Bernstein 型定理. 注意 Enneper 曲面说明嵌入性是不可少的.

我们不证明定理~\ref{thm:2.16}, 而是对嵌入极小圆盘中的内蕴球证明 Colding 与
Minicozzi 在 \cite{CM4} 中得到的一个更强的结果. 这个关于内蕴球的更强估计对
应用至关重要, 并立即给出定理~\ref{thm:2.16} 作为推论. 我们只对 $|A|$ 的界
陈述它, 但注意由推论~\ref{cor:2.7}, 它对伸缩不变的平均面积界也成立.

\begin{theorem}[Colding--Minicozzi, \cite{CM4}]
\label{thm:2.17}
给定常数 $C_I$, 存在 $C_P$, 使得: 若
$B_{2s}^{\Sigma}\subset\Sigma\subset\R^3$ 是嵌入极小圆盘, 且
\begin{equation}
\int_{B_{2s}^{\Sigma}}|A|^{2}\leq C_I,
\label{eq:2.87}
\end{equation}
则
\[
\sup_{B_s^{\Sigma}}|A|^{2}\leq C_Ps^{-2}.
\]
\end{theorem}

该定理容易由下一个引理得到; 该引理对某个环域内总曲率很小的嵌入极小圆盘中的
内蕴球给出曲率估计.

\begin{lemma}[Colding--Minicozzi, \cite{CM4}]
\label{lem:2.18}
给定 $C$, 存在 $\epsilon>0$, 使得: 若
$B_{9s}^{\Sigma}\subset\Sigma\subset\R^3$ 是嵌入极小圆盘, 且
\begin{equation}
\int_{B_{9s}^{\Sigma}}|A|^{2}\leq C
 \quad\text{且}\quad
 \int_{B_{9s}^{\Sigma}\setminus B_s^{\Sigma}}|A|^{2}\leq\epsilon,
\label{eq:2.88}
\end{equation}
则
\[
\sup_{B_s^{\Sigma}}|A|^{2}\leq s^{-2}.
\]
\end{lemma}

\begin{proof}
首先注意, 当 $\epsilon$ 很小时, Choi--Schoen 曲率估计(定理~\ref{thm:2.2})与
\eqref{eq:2.88} 给出
\begin{equation}
\sup_{B_{8s}^{\Sigma}\setminus B_{2s}^{\Sigma}}|A|^{2}
 \leq C_1^{2}\epsilon s^{-2}.
\label{eq:2.89}
\end{equation}
结合 \eqref{eq:2.55} 与 \eqref{eq:2.88} 中第一个不等式, 得到
\begin{equation}
\Length(\partial B_{2s}^{\Sigma})\leq(4\pi+C)s.
\label{eq:2.90}
\end{equation}
下面我们用 \eqref{eq:2.89} 与 \eqref{eq:2.90} 证明: 在旋转 $\R^3$ 之后,
$B_{8s}^{\Sigma}\setminus B_{2s}^{\Sigma}$ (局部地)是 $\{x_3=0\}$ 上的图,
并且
\begin{equation}
\bigl|\Pi(\partial B_{8s}^{\Sigma})\bigr|>3s.
\label{eq:2.91}
\end{equation}
把这两点与嵌入性结合起来, 引理随后容易由 Rado 定理得到.

由 \eqref{eq:2.90}, 我们有直径界
\[
\diam\bigl(B_{8s}^{\Sigma}\setminus B_{2s}^{\Sigma}\bigr)
 \leq(12+2\pi+C/2)s.
\]
于是, 由于
\[
\bigl|\nabla\dist_{\bS^{2}}\bigl(N(x),N\bigr)\bigr|\leq|A|,
\]
对 \eqref{eq:2.89} 积分得到
\begin{equation}
\sup_{x,x'\in B_{8s}^{\Sigma}\setminus B_{2s}^{\Sigma}}
 \dist_{\bS^{2}}\bigl(N(x'),N(x)\bigr)
 \leq C_1\epsilon^{1/2}(12+2\pi+C/2).
\label{eq:2.92}
\end{equation}
因此我们可以旋转 $\R^3$, 使得
\begin{equation}
\sup_{B_{8s}^{\Sigma}\setminus B_{2s}^{\Sigma}}|\nabla x_3|
 \leq C_2\epsilon^{1/2}(1+C).
\label{eq:2.93}
\end{equation}

给定 $y\in\partial B_{2s}^{\Sigma}$, 设 $\gamma_y$ 是从 $y$ 出发、到
$\partial B_{8s}^{\Sigma}$ 的外向法向测地线, 以弧长参数化在 $[0,6s]$ 上.
沿 $\gamma_y$ 对 \eqref{eq:2.89} 积分得到
\begin{equation}
\int_{\gamma_y|_{[0,t]}}\bigl|k_g^{\R^3}\bigr|
 \leq\int_{\gamma_y|_{[0,t]}}|A|
 \leq C_1\epsilon^{1/2}t/s,
\label{eq:2.94}
\end{equation}
其中 $k_g^{\R^3}$ 是 $\gamma_y$ 作为 $\R^3$ 中曲线的测地曲率. 由于
\[
\Bigl|\frac{d}{dt}\langle\gamma_y'(t),\gamma_y'(0)\rangle\Bigr|
 \leq\bigl|k_g^{\R^3}\bigr|,
\]
对 \eqref{eq:2.94} 积分便得到单位法向 $\gamma_y'$ 的振幅界
\[
\langle\gamma_y'(t),\gamma_y'(0)\rangle>1-C_3\epsilon^{1/2}.
\]
再积分一次便得到: $\gamma_y$ 的端点投影到 $\gamma_y'(0)$ 方向直线上相距很远
\begin{equation}
\langle\gamma_y(6s)-\gamma_y(0),\gamma_y'(0)\rangle
 >(1-C_3\epsilon^{1/2})6s.
\label{eq:2.95}
\end{equation}
由于 $\gamma_y(0)\in B_{2s}^{\Sigma}\subset B_{2s}$, 且 \eqref{eq:2.93} 蕴含
$\gamma_y'(0)$ 几乎是水平的, 可见当 $\epsilon$ 小时, \eqref{eq:2.95} 蕴含
\[
\bigl|\Pi(\partial B_{8s}^{\Sigma})\bigr|>3s.
\]

下一步见\figref{fig:2.2}. 结合 $|\Pi(\partial B_{8s}^{\Sigma})|>3s$ 与
\eqref{eq:2.93} 可知, 当 $\epsilon$ 小时,
\[
\Pi^{-1}(\partial D_{2s})\cap B_{8s}^{\Sigma}
\]
是 $\partial D_{2s}$ 上一族浸入的多值图. 由于 $B_{8s}^{\Sigma}$ 是嵌入的,
$\Pi^{-1}(\partial D_{2s})\cap B_{8s}^{\Sigma}$ 由互不相交的嵌入圆周组成, 它们
都是 $\partial D_{2s}$ 上的图; 这是嵌入性唯一的用处. 由于 $x_1^{2}+x_2^{2}$ 在
圆盘 $B_{8s}^{\Sigma}$ 上次调和, 这些圆周在 $B_{8s}^{\Sigma}$ 中张成圆盘, 而由
Rado 定理(即定理~\ref{thm:1.29}), 这些圆盘都是图.

$\partial D_{2s}$ 上这些图的曲率界与梯度界容易推出总测地曲率接近 $2\pi$, 于是
Gauss--Bonnet 给出张成圆盘上的总曲率很小(且当 $\epsilon\to0$ 时趋于零). 由
定理~\ref{thm:2.2} 即得引理.
\end{proof}

\mtfigurepair{fig-2-1}
 {引理~\ref{lem:2.18} 的证明: 由 \eqref{eq:2.93} 与 \eqref{eq:2.94}, 每个 $\gamma_y$ 都几乎是
  长为 $6s$ 的水平线段, 因此
  $|\Pi(\partial B_{8s}^{\Sigma})|>3s$.}{fig:2.1}
 {fig-2-2}
 {引理~\ref{lem:2.18} 的证明:
  $\Pi^{-1}(\partial D_{2s})\cap B_{8s}^{\Sigma}$ 是 $\partial D_{2s}$ 上一族
  图之并, 由 Rado 定理, 每个这样的图在 $\Sigma$ 中张成 $D_{2s}$ 上的一个图.}
 {fig:2.2}

\begin{proof}[定理~\ref{thm:2.17} 的证明]
设 $\epsilon>0$ 由引理~\ref{lem:2.18} 取 $C=C_I$ 给出, 再设 $N$ 是大于
$C_I/\epsilon$ 的最小整数. 给定 $x\in B_s^{\Sigma}$, 存在 $1\leq j\leq N$ 使得
\begin{equation}
\int_{B_{9^{1-j}s}^{\Sigma}(x)\setminus B_{9^{-j}s}^{\Sigma}(x)}|A|^{2}
 \leq\frac{C_I}{N}\leq\epsilon.
\label{eq:2.96}
\end{equation}
结合 \eqref{eq:2.87} 与 \eqref{eq:2.96}, 引理~\ref{lem:2.18} 给出
\[
|A|^{2}(x)\leq(9^{-j}s)^{-2}\leq9^{2N}s^{-2}.
\]
\end{proof}

看起来这一点应当可以推广到高维, 不过显然需要假设比单连通更多的东西. 我们
猜测:

\begin{conjecture}
\label{conj:2.19}
$\R^4$ 中完全、嵌入、可缩且体积三次增长的极小超曲面必定是平坦的.
\end{conjecture}

这一猜测当然可由下面这个更强的猜测推出:

\begin{conjecture}
\label{conj:2.20}
对 $\R^4$ 中可缩且体积有界的嵌入极小超曲面中的内蕴球, 证明一个曲率估计.
(即: 推广定理~\ref{thm:2.17}.)
\end{conjecture}

% ================================================================
\section{稳定超曲面的 $|A|^{2}$ 的 $L^{p}$ 界}
\label{sec:2.4}

本节中, 我们把 Simons 不等式与稳定性不等式结合起来, 对稳定极小超曲面的
第二基本形式模长平方证明更高的 $L^{p}$ 界. 下一节中, 我们将把这个界与 Simons
不等式一起用来证明稳定极小超曲面的曲率估计.

\begin{theorem}[Schoen--Simon--Yau, \cite{ScSiYa}]
\label{thm:2.21}
设 $\Sigma^{n-1}\subset\R^n$ 是可定向稳定极小超曲面. 则对一切
$p\in\bigl[2,2+\sqrt{2/(n-1)}\bigr)$ 及每个非负紧支集 Lipschitz 函数 $\phi$,
\begin{equation}
\int_\Sigma|A|^{2p}\phi^{2p}\leq C(n,p)\int_\Sigma|\nabla\phi|^{2p}.
\label{eq:2.97}
\end{equation}
\end{theorem}

\begin{proof}
若在稳定性不等式(即引理~\ref{lem:1.32})中代入
$\eta=|A|^{1+q}f$($0\leq q<\sqrt{2/(n-1)}$), 得到
\begin{equation}
\begin{aligned}
\int|A|^{4+2q}f^{2}
&\leq\int\bigl|f\nabla|A|^{1+q}+|A|^{1+q}\nabla f\bigr|^{2}\\
&=(1+q)^{2}\int f^{2}\bigl|\nabla|A|\bigr|^{2}|A|^{2q}
 +\int|A|^{2+2q}|\nabla f|^{2}\\
&\quad +2(1+q)\int f|A|^{1+2q}\langle\nabla f,\nabla|A|\rangle.
\end{aligned}
\label{eq:2.98}
\end{equation}
我们将两次用到这个不等式. 关键的一步是 \eqref{eq:2.98} 的第一次使用: 此时把它与 Simons
不等式结合起来, 以估计右端含 $|\nabla|A||^{2}$ 的项. 然后把这一估计代回
\eqref{eq:2.98}, 以证明定理.

把 Simons 不等式
\begin{equation}
|A|\Delta|A|+|A|^{4}\geq\frac{2}{n-1}\bigl|\nabla|A|\bigr|^{2}
\label{eq:2.99}
\end{equation}
乘以 $|A|^{2q}f^{2}$ 并积分, 得到
\begin{equation}
\begin{aligned}
\frac{2}{n-1}\int\bigl|\nabla|A|\bigr|^{2}|A|^{2q}f^{2}
&\leq\int|A|^{4+2q}f^{2}+\int f^{2}|A|^{1+2q}\Delta|A|\\
&=\int|A|^{4+2q}f^{2}
 -2\int f|A|^{1+2q}\langle\nabla f,\nabla|A|\rangle\\
&\quad -(1+2q)\int f^{2}|A|^{2q}\bigl|\nabla|A|\bigr|^{2}.
\end{aligned}
\label{eq:2.100}
\end{equation}
把 \eqref{eq:2.98} 与 \eqref{eq:2.100} 结合起来, 得到
\begin{align}
\Bigl(\frac{2}{n-1}-q^{2}\Bigr)\int|A|^{2q}\bigl|\nabla|A|\bigr|^{2}f^{2}
&\leq\int|A|^{2+2q}|\nabla f|^{2}
 +2q\int f|A|^{1+2q}\langle\nabla f,\nabla|A|\rangle
\label{eq:2.101}\\
&\quad \int|A|^{2+2q}|\nabla f|^{2}
 +\epsilon q\int f^{2}|A|^{2q}\bigl|\nabla|A|\bigr|^{2}
 +\frac{q}{\epsilon}\int|\nabla f|^{2}|A|^{2+2q},
\label{eq:2.102}
\end{align}
其中最后一个不等号用到了吸收不等式 $2xy\leq\epsilon x^{2}+y^{2}/\epsilon$. 于是
得到
\begin{equation}
\Bigl(\frac{2}{n-1}-q^{2}-\epsilon q\Bigr)
 \int f^{2}|A|^{2q}\bigl|\nabla|A|\bigr|^{2}
 \leq\Bigl(1+\frac{q}{\epsilon}\Bigr)\int|\nabla f|^{2}|A|^{2+2q}.
\label{eq:2.103}
\end{equation}
这就完成了第一步.

对 \eqref{eq:2.98} 中的交叉项应用 Cauchy--Schwarz 不等式把它吸收掉, 再代入 \eqref{eq:2.103},
便得到(对任意 $\epsilon<\frac{2/(n-1)-q^{2}}{q}$)
\begin{equation}
\begin{aligned}
\int|A|^{4+2q}f^{2}
&\leq2(1+q)^{2}\int f^{2}|A|^{2q}\bigl|\nabla|A|\bigr|^{2}
 +2\int|A|^{2q+2}|\nabla f|^{2}\\
&\leq\Bigl(\frac{2(1+q)^{2}(1+q/\epsilon)}
 {\frac{2}{n-1}-q^{2}-\epsilon q}+2\Bigr)\int|A|^{2+2q}|\nabla f|^{2}.
\end{aligned}
\label{eq:2.104}
\end{equation}

若置 $p=2+q$, $f=\phi^{p}$, 则 $2\leq p<2+\sqrt{2/(n-1)}$, 而 Hölder
不等式给出
\begin{equation}
\begin{aligned}
\int|A|^{2p}\phi^{2p}
&\leq c\int|A|^{2p-2}\phi^{2p-2}|\nabla\phi|^{2}\\
&\leq c\Bigl(\int|A|^{2p}\phi^{2p}\Bigr)^{\frac{p-1}{p}}
 \Bigl(\int|\nabla\phi|^{2p}\Bigr)^{\frac1p},
\end{aligned}
\label{eq:2.105}
\end{equation}
其中 $c=c(n,p)$.
\end{proof}

% ================================================================
\section{Bernstein 定理与曲率估计}
\label{sec:2.5}

本节中, 我们首先给出第 1 章所证 Bernstein 定理的一个推广.

\begin{theorem}[Schoen--Simon--Yau, \cite{ScSiYa}]
\label{thm:2.22}
若 $\Sigma^{n-1}\subset\R^n$ 是完全双边稳定极小超曲面, $n\leq6$, 且存在
$V<\infty$ 使得
\begin{equation}
\sup_{R>0}\frac{\Vol(B_R\cap\Sigma)}{R^{n-1}}\leq V,
\label{eq:2.106}
\end{equation}
则 $\Sigma$ 是平坦的.
\end{theorem}

\begin{proof}
对每个 $r>0$, 设 $\phi$ 是满足 $\phi|_{B_r}\equiv1$、$\phi|_{B_{2r}}\equiv0$ 的
截断函数, 且在环域 $B_{2r}\setminus B_r$ 上沿径向线性衰减. 把定理~\ref{thm:2.21}
中 $|A|^{2}$ 的 $L^{p}$ 界用于这个截断函数, 其中
\begin{equation}
2p=4+\sqrt{7/5}<4+\sqrt{8/(n-1)},
\label{eq:2.107}
\end{equation}
再与体积界 \eqref{eq:2.106} 结合, 得到
\begin{equation}
\begin{aligned}
\int_{B_r\cap\Sigma}|A|^{4+\sqrt{7/5}}
&\leq C(n,p)\,r^{-4-\sqrt{7/5}}\Vol(B_{2r}\cap\Sigma)\\
&\leq C(n,p)\,2^{n-1}V\,r^{n-5-\sqrt{7/5}}.
\end{aligned}
\label{eq:2.108}
\end{equation}
由于 $n-5-\sqrt{7/5}<0$, 令 $r\to\infty$ 便得到 $|A|^{2}\equiv0$. 命题随即
容易得出.
\end{proof}

由于极小图是稳定的, 且由 \eqref{eq:1.22} 满足面积界 \eqref{eq:2.106}, 我们得到下述 Bernstein
定理:

\begin{theorem}[S. Bernstein \cite{Be}(当 $n=3$), E. De Giorgi \cite{DG}
(当 $n=4$), F. J. Almgren, Jr. \cite{Am1}(当 $n=5$), 以及 J. Simons
\cite{Sim}(当 $n\leq8$)]
\label{thm:2.23}
若 $u:\R^{n-1}\to\R$ 是极小曲面方程的整解, 且 $n\leq8$, 则 $u$ 是线性函数.
\end{theorem}

\begin{proof}
当 $n\leq6$ 时, 这立即由定理~\ref{thm:2.22} 得到. $n=7,8$ 的情形见 \cite{Sim}.
\end{proof}

事实上, 对这一论证稍作变化, 便得到面积有界的稳定极小曲面以及极小图的曲率
估计. 我们只对图给出陈述; 稳定情形的陈述与证明都类似, 只是必须假设面积界.
即: 我们有

\begin{theorem}[E. Heinz \cite{He}(当 $n=3$), Schoen--Simon--Yau \cite{ScSiYa}
(当 $n\leq6$), 以及 L. Simon \cite{Si1}(当 $n\leq8$)]
\label{thm:2.24}
若 $u:D_{r_0}\subset\R^{n-1}\to\R$ 是半径为 $r_0$ 的 $(n-1)$ 维圆盘上极小
曲面方程的解, 且 $n\leq8$, 则对 $0<\sigma\leq r_0$,
\begin{equation}
\sup_{D_{r_0-\sigma}}|A|^{2}\leq C(n)\sigma^{-2}.
\label{eq:2.109}
\end{equation}
\end{theorem}

\begin{proof}
(我们同样只对 $n\leq6$ 证明.) 证明与定理~\ref{thm:2.3} 的证明十分相似. 即: 在
定理~\ref{thm:2.21} 中取 $2p=n$, 并让 $\phi$ 是对数截断函数(如推论~\ref{cor:1.20}
的证明中那样), 便得到: 极小图的面积界蕴含
\begin{equation}
\int_{B_{\sqrt{\kappa}R}\cap\Graph_u}|A|^{n}
 \leq\frac{C_1}{\log\kappa},
\label{eq:2.110}
\end{equation}
其中 $C_1=C_1(n)$, $\kappa>1$, 且 $\kappa R\leq r_0$. 另一方面, Simons 不等式
蕴含
\begin{equation}
\Delta|A|^{n}\geq-C_2|A|^{n+2},
\label{eq:2.111}
\end{equation}
于是我们可以像定理~\ref{thm:2.2} 那样论证, 得到相应的曲率估计.
\end{proof}

% ================================================================
\section{一般的极小图方程}
\label{sec:2.6}

本节的主要结果命题~\ref{prop:2.25} 给出极小曲面 $\Sigma\subset\R^3$ 稳定性的
一个有用判别法. 粗略的想法是: 若一个极小曲面可以从一侧被一列极小曲面逼近,
则它必是稳定的. 这一结果的定量版本在 \cite{CM4} 中分析嵌入极小曲面的局部
结构时起了重要作用.

为把这一点说清楚, 对曲面 $\Sigma$ 上函数 $u$ 的图, 定义
\begin{equation}
\bigl\{x+u(x)N(x)\mid x\in\Sigma\bigr\},
\label{eq:2.112}
\end{equation}
其中 $N$ 是 $\Sigma$ 的单位法向量. 当 $\Sigma$ 是平面时, 这就是通常的图的
概念.

\begin{proposition}
\label{prop:2.25}
设 $\Sigma$ 是极小曲面. 若存在序列 $u_j$, 使得 $u_j$ 在 $\Sigma$ 上的图都是
极小的, $u_j$ 处处不为零, 且 $|u_j|+|\nabla u_j|\to0$, 则 $\Sigma$ 是稳定的.
\end{proposition}

证明该命题的关键在于: 函数 $u_j$ 在 $\Sigma$ 上近似满足 Jacobi 方程.

\begin{lemma}
\label{lem:2.26}
存在常数 $C$, 使得: 若 $\Sigma$ 是极小曲面, $u$ 在 $\Sigma$ 上的图也极小, 且
$\max\{|u||A|,|\nabla u|\}\leq1$, 则 $u$ 满足方程
\begin{equation}
\dive\bigl[(I+\bar L)\nabla u\bigr]+(1+Q)|A|^{2}u+Q_{ij}A_{ij}=0,
\label{eq:2.113}
\end{equation}
其中 $\bar L\leq C(|u||A|+|\nabla u|)$,
$|Q|\leq C(|u||A|+|\nabla u|)^{2}$, 且
$|Q_{ij}|\leq C(|u||A|+|\nabla u|)^{2}$.
\end{lemma}

\begin{proof}
设 $x$ 是位置向量, $N$ 是 $\Sigma$ 的单位法向量, $\{e_1,e_2\}$ 是沿 $\Sigma$ 的
标准正交标架. 我们把 $u$、$e_1$、$e_2$ 与 $N$ 延拓到 $\Sigma$ 的一个小邻域上,
使它们沿法向都为常值. 于是 $\nabla u$ 与 $\Sigma$ 相切, 且
\begin{equation}
\nabla_Ne_i=\nabla_NN=0.
\label{eq:2.114}
\end{equation}
置 $A_{ij}=\langle A(e_i,e_j),N\rangle$, 于是在 $\Sigma$ 的任一点有
\begin{equation}
\nabla_{e_i}N=-A_{ij}e_j.
\label{eq:2.115}
\end{equation}
利用 $\Sigma$ 的极小性与 $A$ 的对称性, 容易看出
\begin{equation}
A_{ij}A_{jk}=\frac12|A|^{2}\delta_{ik}.
\label{eq:2.116}
\end{equation}

$u$ 在 $\Sigma$ 上的图由
\begin{equation}
x\mapsto X(x)=x+u(x)N
\label{eq:2.117}
\end{equation}
给出, 故其切空间由 $X_1$ 与 $X_2$ 张成, 其中
\begin{equation}
X_i=\nabla_{e_i}X=e_i+u_iN-uA_{ik}e_k.
\label{eq:2.118}
\end{equation}
这里我们用 \eqref{eq:2.115} 对 $N$ 求导, 并用到对任意向量 $v$ 都有 $\nabla_vx=v$.

利用 \eqref{eq:2.118}, 我们算出该图的度量 $g_{ij}$:
\begin{equation}
\begin{aligned}
g_{ij}&=\langle X_i,X_j\rangle
 =\delta_{ij}+u_iu_j+u^{2}A_{ik}A_{jk}-2uA_{ij}\\
&=\Bigl(1+\frac{|A|^{2}}{2}u^{2}\Bigr)\delta_{ij}+u_iu_j-2uA_{ij}.
\end{aligned}
\label{eq:2.119}
\end{equation}

由于我们要考虑的是 $|\nabla u|+|u||A|$ 很小的解, 我们将经常把它写成
\begin{equation}
g_{ij}=\delta_{ij}-2uA_{ij}+Q_{ij},
\label{eq:2.120}
\end{equation}
其中 $Q_{ij}$ 表示一个关于 $|\nabla u|$ 与 $|u||A|$ 二次的矩阵(并且允许它在
不同行中有所不同). 更确切地说, 存在 $C$, 使得当
$\max\{|u||A|,|\nabla u|\}\leq1$ 时,
\begin{equation}
\bigl|Q_{ij}(x)\bigr|\leq C\bigl(|u||A|+|\nabla u|\bigr)^{2}.
\label{eq:2.121}
\end{equation}
由此可见逆度量为
\begin{equation}
g^{ij}=\delta_{ij}+2uA_{ij}+Q_{ij}.
\label{eq:2.122}
\end{equation}

由 \eqref{eq:2.120} 可得
\begin{equation}
\det g_{ij}=(1-2uA_{11})(1-2uA_{22})+Q=1+Q,
\label{eq:2.123}
\end{equation}
其中第二个等号用到了极小性(即 $A_{11}+A_{22}=0$), 而这次的 $Q$ 是一个标量二次
项. 若 $J(s)$ 是一族可微的矩阵, 则行列式在 $0$ 处的导数为
\begin{equation}
\frac{d}{ds}\Big|_{s=0}\det J(s)
 =\det J(0)\operatorname{Trace}\bigl(J^{-1}(0)J'(0)\bigr).
\label{eq:2.124}
\end{equation}
把它用于 $J(s)=g_{ij}(s)$, 其中 $g_{ij}(s)$ 是把 $u$ 换成 $u+sv$ 后算出的, 我们
得到
\begin{equation}
\begin{aligned}
\frac{d}{ds}\Big|_{s=0}&\det g_{ij}\\
&=(1+Q)(\delta_{ij}+2uA_{ij}+Q_{ij})
 \bigl(|A|^{2}uv\delta_{ij}-2vA_{ij}+u_iv_j+u_jv_i\bigr)\\
&=2\langle\nabla u,\nabla v\rangle-2(1+Q)|A|^{2}uv
 +Q_{ij}A_{ij}v+\langle\bar L\nabla u,\nabla v\rangle,
\end{aligned}
\label{eq:2.125}
\end{equation}
其中最后一项里的 $\bar L$ 是一个关于 $|\nabla u|+|A||u|$ 线性有界的矩阵(仍假设
$\max\{|u||A|,|\nabla u|\}\leq1$), 并且我们用到了 $(1+Q)Q_{ij}=Q_{ij}$. 因此,
若 $u$ 的图是极小的, 则对每个紧支集 $v$ 都得到
\begin{equation}
0=\int_\Sigma\langle(I+\bar L)\nabla u,\nabla v\rangle
 -(1+Q)|A|^{2}uv+Q_{ij}A_{ij}v.
\label{eq:2.126}
\end{equation}
分部积分把 $\nabla$ 从 $v$ 上移开, 便得 $u$ 满足 \eqref{eq:2.113}.
\end{proof}

\begin{proof}[命题~\ref{prop:2.25} 的证明]
由命题~\ref{prop:1.39}, 只需在任意一个闭包紧致的连通开子集 $K_0\subset\Sigma$
上构造一个处处不为零的 Jacobi 场 $w$. 固定 $K_0$ 内部的一点 $p$, 并设 $K_1$
是包含 $K_0$ 的连通开集, 其闭包也紧致.

由于 $|A|$ 在 $\Sigma$ 的紧子集上有界, 且
\begin{equation}
|u_j|+|\nabla u_j|\to0,
\label{eq:2.127}
\end{equation}
引理~\ref{lem:2.26} 给出 $C$, 使得 $u_j$ 满足
\begin{equation}
\dive\bigl[(I+\bar L_j)\nabla u_j\bigr]
 +(1+Q_j)|A|^{2}u_j+\bar Q_j=0,
\label{eq:2.128}
\end{equation}
其中
\begin{equation}
|\bar L_j|^{2}+|Q_j|+|\bar Q_j|
 \leq C\bigl(|u||A|+|\nabla u|\bigr)^{2}.
\label{eq:2.129}
\end{equation}
特别地, $u_j$ 满足统一的 Harnack 不等式:
\begin{equation}
\sup_{K_1}u_j\leq C_1\inf_{K_1}u_j\leq C_1u_j(p).
\label{eq:2.130}
\end{equation}
在二维情形, 上述方程还蕴含 $u_j$ 的梯度的 Hölder 连续性(见 \cite{GiTr} 定理
12.4), 从而使我们可以应用 Schauder 估计(\cite{GiTr} 定理 6.2)得到
\begin{equation}
\|u_j\|_{C^{2,\alpha}(K_0)}\leq C_2u_j(p).
\label{eq:2.131}
\end{equation}
特别地, 若定义函数 $w_j$ 为
\begin{equation}
w_j(x)=\frac{u_j(x)}{u_j(p)},
\label{eq:2.132}
\end{equation}
则每个 $w_j(p)=1$, $w_j$ 都是正的, 且
\begin{equation}
\|w_j\|_{C^{2,\alpha}(K_0)}\leq C_2.
\label{eq:2.133}
\end{equation}
因此我们可以应用 Arzelà--Ascoli 定理, 取出一子列(仍记作 $w_j$), 它在 $C^{2}$
中一致收敛于极限函数 $w$, 且 $w(p)=1$, $w\geq0$. 把 \eqref{eq:2.128} 乘以 $1/u_j(p)$
并利用一致收敛性, 可以取极限得到
\begin{equation}
\dive\nabla w+|A|^{2}w=0,
\label{eq:2.134}
\end{equation}
其中我们用到了 \eqref{eq:2.129} 以及 $|u_j|+|\nabla u_j|\to0$, 以消去极限中的高阶项.
最后, $w$ 也满足 Harnack 不等式, 故它必定也是正的.
\end{proof}

这些结果也都有定量版本, 即可以证明: 若两个曲率有界的极小曲面靠得很近但不相交,
则它们各自都几乎是稳定的. 下面我们来讨论这一点.

% ================================================================
\section{几乎稳定}
\label{sec:2.7}

稳定极小曲面的估计可以推广到在 \cite{CM4} 的意义下``几乎稳定''的曲面, 其确切
含义下面给出. 这一推广至关重要, 因为可以证明: 曲率有界的、彼此接近但不相交的
极小曲面是几乎稳定的, 却未必按通常意义稳定. 本节结果取自 \cite{CM4}.

我们先给出稳定性的一个推广:

\begin{definition}[$\delta_s$-稳定]
\label{def:2.27}
给定 $\delta_s\geq0$, 置
\begin{equation}
L_{\delta_s}=\Delta+(1-\delta_s)|A|^{2},
\label{eq:2.135}
\end{equation}
于是 $L_0$ 就是 $\Sigma$ 上通常的 Jacobi 算子. 若对任意紧支集 Lipschitz 函数
$\phi$(即 $\phi\in C_0^{0,1}(\Omega)$)都有
\[
\int\phi L_{\delta_s}\phi\leq0,
\]
则称区域 $\Omega\subset\Sigma$ 是 \textbf{$\delta_s$-稳定}的.
\end{definition}

由此可见, $\Omega$ 是 $\delta_s$-稳定的, 当且仅当对一切
$\phi\in C_0^{0,1}(\Omega)$ 都有 $\delta_s$-稳定性不等式:
\begin{equation}
(1-\delta_s)\int|A|^{2}\phi^{2}\leq\int|\nabla\phi|^{2}.
\label{eq:2.136}
\end{equation}

由于由引理~\ref{lem:2.26}, Jacobi 方程就是 $\Sigma$ 上极小图方程的线性化,
标准的计算给出:

\begin{lemma}
\label{lem:2.28}
存在 $\delta_g>0$, 使得: 若 $\Sigma$ 极小, $u$ 是 $\Sigma$ 上极小图方程的正解
(即 $\{x+u(x)N_{\Sigma}(x)\mid x\in\Sigma\}$ 是极小的), 且
\[
|\nabla u|+|u||A|\leq\delta_g,
\]
则 $w=\log u$ 在 $\Sigma$ 上满足
\begin{equation}
\Delta w=-|\nabla w|^{2}+\dive(a\nabla w)+\langle\nabla w,a\nabla w\rangle
 +\langle b,\nabla w\rangle+(c-1)|A|^{2},
\label{eq:2.137}
\end{equation}
其中 $a_{ij}$、$b_j$、$c$ 是 $\Sigma$ 上的函数, 满足
$|a|,|c|\leq3|A||u|+|\nabla u|$, $|b|\leq2|A||\nabla u|$.
\end{lemma}

\begin{proof}
这容易由引理~\ref{lem:2.26} 得到.
\end{proof}

下面这个对命题~\ref{prop:1.39} 的微改给出区域 $\delta_s$-稳定性的一个有用的
充分条件:

\begin{lemma}[Colding--Minicozzi, \cite{CM4}]
\label{lem:2.29}
存在 $\delta>0$, 使得: 若 $\Sigma$ 极小, $u>0$ 是 $\Omega\subset\Sigma$ 上
极小图方程的解, 且
\[
|\nabla u|+|u||A|\leq\delta,
\]
则 $\Omega$ 是 $1/2$-稳定的.
\end{lemma}

\begin{proof}
置 $w=\log u$, 并取截断函数 $\phi\in C_0^{0,1}(\Omega)$. 对
\[
\dive\bigl(\phi^{2}\nabla w-\phi^{2}a\nabla w\bigr)
\]
应用 Stokes 定理, 代入 \eqref{eq:2.137}, 并利用 $|a|,|c|\leq3\delta$,
$|b|\leq2\delta|\nabla w|$, 得到
\begin{equation}
\begin{aligned}
(1-3\delta)\int\phi^{2}|A|^{2}
&\leq-\int\phi^{2}|\nabla w|^{2}
 +\int\phi^{2}\langle\nabla w,b+a\nabla w\rangle\\
&\quad +2\int\phi\langle\nabla\phi,\nabla w-a\nabla w\rangle\\
&\leq(5\delta-1)\int\phi^{2}|\nabla w|^{2}
 +2(1+3\delta)\int|\phi\nabla w||\nabla\phi|.
\end{aligned}
\label{eq:2.138}
\end{equation}
引理由吸收不等式 $2xy\leq\epsilon x^{2}+y^{2}/\epsilon$ 得到.
\end{proof}

我们将用引理~\ref{lem:2.29} 说明: 彼此接近的、不相交的嵌入极小曲面是近似
稳定的(即下面的推论~\ref{cor:2.31}).

我们还需要把引理~\ref{lem:2.4}$|A|$ 的界给出图形化区域这一结论精细化. 与
引理~\ref{lem:2.4} 的证明一样, 出发点是把
\begin{equation}
\bigl|\nabla\dist_{\bS^{2}}\bigl(N(x),N\bigr)\bigr|\leq|A|
\label{eq:2.139}
\end{equation}
沿测地线积分, 得到
\begin{equation}
\sup_{x'\in B_s^{\Sigma}(x)}\dist_{\bS^{2}}\bigl(N(x'),N(x)\bigr)
 \leq s\sup_{B_s^{\Sigma}(x)}|A|.
\label{eq:2.140}
\end{equation}

由 \eqref{eq:2.140}, 我们可以选取 $0<\rho_2<1/4$, 使得:

若 $B_{2s}^{\Sigma}(x)\subset\Sigma$,
$s\sup_{B_{2s}^{\Sigma}(x)}|A|\leq4\rho_2$, 且 $t\leq s$, 则
$B_t(x)\cap\Sigma$ 中含 $x$ 的连通分支 $\Sigma_{x,t}$ 是 $T_x\Sigma$ 上的图,
其梯度不超过 $t/s$, 并且
\begin{equation}
\inf_{x'\in B_{2s}^{\Sigma}(x)}
 \frac{|x'-x|}{\dist_{\Sigma}(x,x')}>\frac{9}{10}.
\label{eq:2.141}
\end{equation}

其一个推论是: 若 $t\leq s$, 并把 $T_x\Sigma$ 平移使 $x\in T_x\Sigma$, 则
\begin{equation}
\sup_{x'\in B_t^{\Sigma}(x)}|x'-T_x\Sigma|\leq t^{2}/s.
\label{eq:2.142}
\end{equation}

下一个引理说明: 两个曲率有界、彼此靠近的极小曲面可以互相写成对方的图, 且这些
函数近似满足 Jacobi 方程. 该引理的表述方式使它既可用于两个处处不相交的曲面,
也可用于单个嵌入曲面的两个子集, 只要这两个子区域在内蕴意义下相距足够远.

\begin{lemma}[Colding--Minicozzi, \cite{CM4}]
\label{lem:2.30}
存在 $C_0,\rho_0>0$, 使得下述结论成立:

若 $\rho_1\leq\min\{\rho_0,\rho_2\}$, $\Sigma_1,\Sigma_2\subset\R^3$ 是定向
极小曲面, 在各自 $\Sigma_i$ 上 $|A|^{2}\leq4$, 并且
\begin{gather*}
x\in\Sigma_1\setminus\tub_{4\rho_2}(\partial\Sigma_1),\\[3pt]
y\in B_{\rho_1}(x)\cap\Sigma_2\setminus\tub_{4\rho_2}(\partial\Sigma_2),\\[3pt]
B_{2\rho_1}^{\Sigma}(x)\cap B_{2\rho_1}^{\Sigma}(y)=\emptyset,
\end{gather*}
则 $B_{\rho_2}^{\Sigma}(y)$ 是 $\{z+u(z)N(z)\}$ 在某个包含
$B_{\rho_2/2}^{\Sigma}(x)$ 的区域上的图, 其中 $u\neq0$ 且
$|\nabla u|+4|u|\leq C_0\rho_1$.
\end{lemma}

\begin{proof}
由于 $\rho_1\leq\rho_2$, \eqref{eq:2.141} 蕴含
\[
B_{2\rho_2}^{\Sigma}(x)\cap B_{2\rho_2}^{\Sigma}(y)=\emptyset.
\]
若 $t\leq9\rho_2/5$, 则 $|A|^{2}\leq4$ 蕴含 $B_t(x)\cap\Sigma_1$ 与
$B_t(y)\cap\Sigma_2$ 中含 $x$、$y$ 的连通分支 $\Sigma_{x,t}$、$\Sigma_{y,t}$
分别是 $T_x\Sigma_1$、$T_y\Sigma_2$ 上的图, 梯度不超过 $t/(2\rho_2)$, 并且
\[
\Sigma_{x,t}\subset B_{2\rho_2}^{\Sigma}(x),\qquad
\Sigma_{y,t}\subset B_{2\rho_2}^{\Sigma}(y).
\]
最后的结论蕴含 $\Sigma_{x,t}\cap\Sigma_{y,t}=\emptyset$. 由此可知
$\Sigma_{x,t}$ 与 $\Sigma_{y,t}$ 是同一平面上的图. 即: 若置
\[
\theta=\dist_{\bS^{2}}\bigl(N(x),\{N(y),-N(y)\}\bigr),
\]
则由 \eqref{eq:2.142}、$|x-y|<\rho_1$ 以及 $\Sigma_{x,t}\cap\Sigma_{y,t}=\emptyset$
蕴含
\begin{equation}
\rho_1-(t/2-\rho_1)\sin\theta+t^{2}/(2\rho_2)>-t^{2}/(2\rho_2).
\label{eq:2.143}
\end{equation}
于是
\[
\sin\theta<\rho_1/(t/2-\rho_1)
 +t^{2}\bigl/[(t/2-\rho_1)\rho_2\bigr].
\]
当 $\rho_0/\rho_2$ 很小时, $B_{\rho_2}^{\Sigma}(y)$ 是 $T_x\Sigma_1$ 上梯度有界
的图. 引理现在容易由 Harnack 不等式得到.
\end{proof}

\mtfigurepair{fig-2-3}
 {推论~\ref{cor:2.31}: 两个足够接近、曲率有界且不相交的极小曲面必定是近似
  稳定的.}{fig:2.3}
 {fig-2-4}
 {\eqref{eq:2.148} 中的集合 $VB$. 这里 $x\in VB$, 而 $y\in\Sigma\setminus VB$.}{fig:2.4}

结合引理~\ref{lem:2.29} 与引理~\ref{lem:2.30}, 得到下一个推论:

\begin{corollary}[Colding--Minicozzi, \cite{CM4}]
\label{cor:2.31}
见\figref{fig:2.3}. 给定 $C_0$、$\delta>0$, 存在
$\epsilon(C_0,\delta)>0$, 使得: 若 $p_i\in\Sigma_i\subset\R^3$($i=1,2$)是
嵌入极小曲面, 满足
\[
\Sigma_1\cap\Sigma_2=\emptyset,\quad
B_{2R}^{\Sigma}(p_i)\cap\partial\Sigma_i=\emptyset,\quad
|p_1-p_2|<\epsilon R
\]
以及
\begin{equation}
\sup_{B_{2R}^{\Sigma}(p_i)}|A|^{2}\leq C_0R^{-2},
\label{eq:2.144}
\end{equation}
则 $B_R^{\Sigma}(\tilde p_i)\subset\tilde\Sigma_i$ 是 $\delta$-稳定的, 其中
$\tilde p_i$ 是 $p_i$ 在 $\Sigma_i$ 的万有覆盖 $\tilde\Sigma_i$ 中的原像点.
\end{corollary}

\mtpara{寻找大的几乎稳定区域}
下一个结果把曲率有界的嵌入极小曲面分解为两部分: 面积有界的一部分, 以及一族
互不相交的 $1/2$-稳定区域之并. 这个结果技术性更强, 但在后面几章研究嵌入极小
圆盘在大曲率点附近的局部结构(如 \cite{CM4})时会用到.

\begin{lemma}[Colding--Minicozzi, \cite{CM4}]
\label{lem:2.32}
存在 $C_1$, 使得:

若 $0\in\Sigma\subset B_{2R}\subset\R^3$ 是嵌入极小曲面,
$\partial\Sigma\subset\partial B_{2R}$, 且 $|A|^{2}\leq4$, 则存在互不相交的
$1/2$-稳定子区域 $\Omega_j\subset\Sigma$ 以及函数 $\chi\leq1$, 使得 $\chi$ 在
$B_R\cap\Sigma\setminus\cup_j\Omega_j$ 上为零, 并且
\begin{equation}
\Area\bigl(\{x\in B_R\cap\Sigma\mid\chi(x)<1\}\bigr)\leq C_1R^{3},
\label{eq:2.145}
\end{equation}
\begin{equation}
\int_{B_R^{\Sigma}}|\nabla\chi|^{2}\leq C_1R^{3}.
\label{eq:2.146}
\end{equation}
\end{lemma}

\begin{proof}
可以假设 $R>\rho_2$(否则 $B_R\cap\Sigma$ 是稳定的). 设 $\delta>0$ 来自
引理~\ref{lem:2.29}, $C_0,\rho_0$ 来自引理~\ref{lem:2.30}. 置
\begin{equation}
\rho_1=\min\Bigl\{\frac{\rho_0}{C_0},\frac{\delta}{C_0},
 \frac{\rho_2}{4}\Bigr\}.
\label{eq:2.147}
\end{equation}
给定 $x\in B_{2R-\rho_1}\cap\Sigma$, 设 $\Sigma_x$ 是 $B_{\rho_1}(x)\cap\Sigma$
中含 $x$ 的连通分支, $B_x^{+}$ 是 $B_{\rho_1}(x)\setminus\Sigma_x$ 中
$N(x)$ 所指入的连通分支. 见\figref{fig:2.4}. 置
\begin{equation}
VB=\bigl\{x\in B_R\cap\Sigma\mid
 B_x^{+}\cap\Sigma\setminus B_{4\rho_1}^{\Sigma}(x)=\emptyset\bigr\},
\label{eq:2.148}
\end{equation}
并设 $\{\Omega_j\}$ 是 $B_R\cap\Sigma\setminus\overline{VB}$ 的连通分支. 选取
一个极大的互不相交球族
\begin{equation}
\bigl\{B_{\rho_1}^{\Sigma}(y_i)\bigr\}_{1\leq i\leq\nu},
\label{eq:2.149}
\end{equation}
球心都在 $VB$ 中. 于是球族
\begin{equation}
\bigl\{B_{2\rho_1}^{\Sigma}(y_i)\bigr\}_{1\leq i\leq\nu}
\label{eq:2.150}
\end{equation}
覆盖集合 $VB$. 此外, ``半球''
\begin{equation}
B_{\rho_1/2}(y_i)\cap B_{y_i}^{+}
\label{eq:2.151}
\end{equation}
两两不相交. 要看出这一点, 假设 $|y_i-y_j|<\rho_1$, 但
$y_j\notin B_{2\rho_1}^{\Sigma}(y_i)$. 则由 \eqref{eq:2.141}, $y_j\notin
B_{8\rho_1}^{\Sigma}(y_i)$, 故 $B_{4\rho_1}^{\Sigma}(y_j)\notin B_{y_i}^{+}$,
$B_{4\rho_1}^{\Sigma}(y_i)\notin B_{y_j}^{+}$; 三角不等式随后蕴含
\[
B_{\rho_1/2}(y_i)\cap B_{y_i}^{+}\cap B_{\rho_1/2}(y_j)\cap B_{y_j}^{+}
 =\emptyset,
\]
如所断言. 由 \eqref{eq:2.140}--\eqref{eq:2.142} 可见, 每个
$B_{\rho_1/2}(y_i)\cap B_{y_i}^{+}$ 的体积约为 $\rho_1^{3}$, 且含于
$B_{2R}$ 中, 于是 $\nu\leq CR^{3}$. 定义 $\Sigma$ 上的函数 $\chi$ 为
\begin{equation}
\chi(x)=\begin{cases}
0 & \text{若 } x\in VB,\\
\dist_{\Sigma}(x,VB)/\rho_1
 & \text{若 } x\in\tub_{\rho_1}(VB)\setminus VB,\\
1 & \text{其他情形}.
\end{cases}
\label{eq:2.152}
\end{equation}
由于 $\tub_{\rho_1}(VB)\subset\bigcup_{i=1}^{\nu}
B_{3\rho_1}^{\Sigma}(y_i)$, $|A|^{2}\leq4$ 且 $\nu\leq CR^{3}$, 我们得到
\eqref{eq:2.145}. 结合 \eqref{eq:2.145} 与 $|\nabla\chi|\leq\rho_1^{-1}$ 便得到 \eqref{eq:2.146}(把 $C_1$
取得更大些).

剩下要证明每个 $\Omega_j$ 是 $1/2$-稳定的. 固定 $j$. 按构造, 若
$x\in\Omega_j$, 则存在 $y_x\in B_x^{+}\cap\Sigma\setminus
B_{4\rho_1}^{\Sigma}(x)$, 它在 $B_x^{+}\cap\Sigma$ 中使 $|x-y_x|$ 取到最小.
特别地, 由引理~\ref{lem:2.30}, $B_{\rho_2}^{\Sigma}(y_x)$ 是
$\{z+u_x(z)N(z)\}$ 在包含 $B_{\rho_2/2}^{\Sigma}(x)$ 的区域上的图, 其中
$u_x>0$ 且
\begin{equation}
|\nabla u_x|+4|u_x|\leq\min\{\delta,\rho_0\}.
\label{eq:2.153}
\end{equation}
选取一个极大的互不相交球族 $B_{\rho_2/6}^{\Sigma}(x_i)$, $x_i\in\Omega_j$, 并
设 $u_{x_i}>0$ 是对应定义在 $B_{\rho_2/2}^{\Sigma}(x_i)$ 上的函数. 由于 $\Sigma$
是嵌入的(且紧), 而 $|u_{x_i}|<\rho_0$, 引理~\ref{lem:2.30} 蕴含
\[
u_{x_i}(x)=\min_{t>0}\{x+tN(x)\in\Sigma\}
\]
对 $x\in B_{\rho_2/2}^{\Sigma}(x_i)$ 成立. 因此, 当
$x\in B_{\rho_2/2}^{\Sigma}(x_{i_1})\cap B_{\rho_2/2}^{\Sigma}(x_{i_2})$ 时
$u_{x_{i_1}}(x)=u_{x_{i_2}}(x)$. 注意
$\tub_{\rho_2/6}(\Omega_j)\subset\bigcup_i
B_{\rho_2/2}^{\Sigma}(x_i)$. 我们断定: 这些 $u_{x_i}$ 给出
$\tub_{\rho_2/6}(\Omega_j)$ 上一个定义良好的正函数 $u_j$, 满足
\begin{equation}
|\nabla u_j|+|u_j||A|\leq\delta.
\label{eq:2.154}
\end{equation}
最后, 引理~\ref{lem:2.29} 蕴含每个 $\Omega_j$ 是 $1/2$-稳定的.
\end{proof}

% ================================================================
\section{间距的次线性增长}
\label{sec:2.8}

如我们所见, 对来自螺旋面两半的多值图, 其间距是常数. 这可以看作一种 Liouville
型定理, 反映了无穷值图的共形性质. 在 \cite{CM3} 中, 我们对定义在多值图上的
极小图方程的正解给出了相应的梯度估计:

\begin{theorem}[Colding--Minicozzi, \cite{CM3}]
\label{thm:2.33}
给定 $\alpha>0$, 存在 $\delta_p>0$ 与 $N_g>5$, 使得下述结论成立:

若 $\Sigma$ 是函数 $u$ 在 $D_{\mathrm{e}^{N_g}R}\setminus
D_{\mathrm{e}^{-N_g}R}$ 上的 $N_g$ 值极小图, $|\nabla u|\leq1$, 且
$0<w<\delta_pR$ 是 $\Sigma$ 上极小图方程的解, $|\nabla w|\leq1$, 则对
$R\leq s\leq2R$,
\begin{equation}
\sup_{\Sigma^{0,2\pi}_{R,2R}}|A_{\Sigma}|
 +\sup_{\Sigma^{0,2\pi}_{R,2R}}|\nabla w|/w
 \leq\alpha/(4R),
\label{eq:2.155}
\end{equation}
\begin{equation}
\sup_{\Sigma^{0,2\pi}_{R,s}}w
 \leq(s/R)^{\alpha}\sup_{\Sigma^{0,2\pi}_{R,R}}w.
\label{eq:2.156}
\end{equation}
\end{theorem}

我们将在本节证明这个定理.

\eqref{eq:2.156} 的一个重要推论是: 当 $\alpha<1$ 时, 间距的增长是次线性的. 这一点在
\cite{CM3} 中推广多值图时起了重要作用, 在 \cite{CM6} 中证明单侧曲率估计时
(见第 8 章)又一次起了作用. 关于层数增长的多值图更精细的对数估计, 见
\cite{CM8}; 并参看 \cite{CM29}. 当多值图是某个嵌入极小圆盘的一部分时, 还可以
说得更多; 见 \cite{CM8}.

设 $\mathcal{P}$ 是去心平面 $\mathbb{C}\setminus\{0\}$ 的万有覆盖, 带整体极坐标
$(\rho,\theta)$, 于是 $\rho>0$, $\theta\in\R$.

\mtpara{调和函数的类比}
在证明定理~\ref{thm:2.33} 之前, 我们在本小节说明它为什么成立.

函数 $w$ 近似是线性化方程(在这里就是 Jacobi 方程)的解. 由于这些图有梯度界,
该方程与 Laplace 方程相距不远. 为说明 \eqref{eq:2.156} 为什么成立, 我们在 $u$ 与 $w$
都调和的条件下给出一个初等证明.

\begin{proof}[$u$ 调和时定理~\ref{thm:2.33} 的证明]
经伸缩, 可设 $R=1$. 作共形坐标变换
\begin{equation}
(\rho,\theta)\mapsto(\log\rho,\theta),
\label{eq:2.157}
\end{equation}
我们得到正方形 $[-N_g,N_g]\times[-N_g,N_g]$ 上的正调和函数
\begin{equation}
\tilde w(x,y)=w\bigl(\mathrm{e}^{x},y\bigr).
\label{eq:2.158}
\end{equation}
由于链式法则给出
\begin{equation}
\nabla\log w(1,0)=\nabla\log\tilde w(0,0),
\label{eq:2.159}
\end{equation}
把欧氏梯度估计用于 $\tilde w$, 得到
\begin{equation}
\bigl|\nabla\log w(1,0)\bigr|
 =\bigl|\nabla\log\tilde w(0,0)\bigr|\leq C/N_g.
\label{eq:2.160}
\end{equation}
若 $N_g$ 足够大, 这就给出 $w$ 的次线性梯度估计. 对 $\Hess_u$ 的界类似可得.
\end{proof}

\mtpara{构造好的截断函数}
下面这个引理与推论构造了定理~\ref{thm:2.33} 的梯度估计所需的截断函数. 本小节
与下一小节的材料都取自 \cite{CM3}.

\begin{lemma}
\label{lem:2.34}
给定 $N>36/(1-\mathrm{e}^{-1/3})^{2}$, 存在 $\mathcal{P}$ 上的函数
$0\leq\phi\leq1$, 满足 $\int|\nabla\phi|^{2}\leq4\pi/\log N$, 且
\begin{equation}
\phi=\begin{cases}
1 & \text{若 } R/\mathrm{e}\leq\rho\leq\mathrm{e}R
 \text{ 且 } |\theta|\leq3\pi,\\
0 & \text{若 } \rho\leq\mathrm{e}^{-N}R
 \text{ 或 } \mathrm{e}^{N}R\leq\rho
 \text{ 或 } |\theta|\geq\pi N.
\end{cases}
\label{eq:2.161}
\end{equation}
\end{lemma}

\begin{proof}
经伸缩, 可设 $R=1$. 由于曲面上的能量在共形变换下不变, 与 $z^{3N}$ 复合可知
\eqref{eq:2.161} 等价于 $E(\phi)\leq4\pi/\log N$ 以及
\begin{equation}
\phi=\begin{cases}
1 & \text{若 } |\log\rho|<1/(3N) \text{ 且 } |\theta|\leq\pi/N,\\
0 & \text{若 } |\log\rho|>1/3 \text{ 或 } |\theta|\geq\pi/3.
\end{cases}
\label{eq:2.162}
\end{equation}
若置
\begin{equation}
\phi=\begin{cases}
1 & \text{于 } D_{6/N}(1,0),\\
1-\dfrac{\log[N\dist_{\mathcal{P}}((1,0),\cdot)/6]}
 {\log[N(1-\mathrm{e}^{-1/3})/6]}
 & \text{于 } D_{1-\mathrm{e}^{-1/3}}(1,0)\setminus D_{6/N}(1,0),\\
0 & \text{其他情形},
\end{cases}
\label{eq:2.163}
\end{equation}
则上述要求就能满足(此时 $E(\phi)=2\pi/\log[N(1-\mathrm{e}^{-1/3})/6]$).
\end{proof}

给定一个 $N$ 值图 $\Sigma$, 用 $\Sigma^{\theta_1,\theta_2}_{r_3,r_4}\subset\Sigma$
表示它在
\begin{equation}
\bigl\{(\rho,\theta)\mid r_3\leq\rho\leq r_4,\ \theta_1\leq\theta
 \leq\theta_2\bigr\}
\label{eq:2.164}
\end{equation}
上的子图.

把引理~\ref{lem:2.34} 的截断函数移植到多值图上, 便得到下一个推论:

\begin{corollary}
\label{cor:2.35}
给定 $\epsilon_0,\tau>0$, 存在 $N>0$, 使得: 若 $\Sigma\subset\R^3$ 是 $u$ 在
$D_{\mathrm{e}^{N}R}\setminus D_{\mathrm{e}^{-N}R}$ 上的 $N$ 值图, 且
$|\nabla u|\leq\tau$, 则 $\Sigma$ 上存在截断函数 $0\leq\phi\leq1$, 满足
$\int|\nabla\phi|^{2}\leq\epsilon_0$, $\phi|_{\partial\Sigma}=0$, 并且
\begin{equation}
\phi\equiv1\quad\text{于 }\ \Sigma^{-\pi,3\pi}_{R/2,5R/2}.
\label{eq:2.165}
\end{equation}
\end{corollary}

\begin{proof}
由于 $\Sigma^{-\pi,3\pi}_{R/2,5R/2}\subset\Sigma^{-3\pi,3\pi}_{R/\mathrm{e},
\mathrm{e}R}$, 且从 $\Sigma$ 到 $\mathcal{P}$ 的投影是双 Lipschitz 的, 其双
Lipschitz 常数不超过 $\sqrt{1+\tau^{2}}$, 故推论由引理~\ref{lem:2.34} 得到.
\end{proof}

\mtpara{梯度估计}
若 $u>0$ 是 $\Sigma$ 上 Jacobi 方程 $\Delta u=-|A|^{2}u$ 的解, 则 $w=\log u$
满足
\begin{equation}
\Delta w=-|\nabla w|^{2}-|A|^{2}.
\label{eq:2.166}
\end{equation}
Bochner 公式(命题~\ref{prop:1.47})、\eqref{eq:2.166}、$\mathrm{K}_{\Sigma}=-|A|^{2}/2$
以及 Cauchy--Schwarz 不等式给出
\begin{equation}
\begin{aligned}
\Delta|\nabla w|^{2}
&=2|\Hess_w|^{2}+2\langle\nabla w,\nabla\Delta w\rangle-|A|^{2}|\nabla w|^{2}\\
&\geq2|\Hess_w|^{2}-4|\nabla w|^{2}|\Hess_w|
 -4|\nabla w||A||\nabla A|-|A|^{2}|\nabla w|^{2}\\
&\geq-2|\nabla w|^{4}-3|A|^{2}|\nabla w|^{2}-2|\nabla A|^{2}.
\end{aligned}
\label{eq:2.167}
\end{equation}
利用 Jacobi 方程是 $\Sigma$ 上极小图方程的线性化这一点, 引理~\ref{lem:2.28}
表明 \eqref{eq:2.166} 与 \eqref{eq:2.167} 的类似式对 $\Sigma$ 上极小图方程的解也成立.

\begin{proof}[定理~\ref{thm:2.33} 的证明]
固定 $\epsilon_E>0$(稍后取为只依赖于 $\alpha$). 推论~\ref{cor:2.35} 给出 $N$
(只依赖于 $\epsilon_E$)以及
$\Sigma^{-N\pi,N\pi}_{\mathrm{e}^{-N}R,\mathrm{e}^{N}R}$ 上紧支集的函数
$0\leq\phi\leq1$, 满足
\begin{equation}
\int|\nabla\phi|^{2}\leq\epsilon_E,
\quad\text{且}\quad
\phi\equiv1\ \text{于 }\ \Sigma^{-\pi,3\pi}_{R/2,5R/2}.
\label{eq:2.168}
\end{equation}
置 $N_g=N+1$, 于是
\[
\dist_{\Sigma}\Bigl(\Sigma^{-N\pi,N\pi}_{\mathrm{e}^{-N}R,\mathrm{e}^{N}R},
 \partial\Sigma\Bigr)>\mathrm{e}^{-N}R/2,
\]
从而 Heinz 曲率估计给出: 在
$\Sigma^{-N\pi,N\pi}_{\mathrm{e}^{-N}R,\mathrm{e}^{N}R}$ 上
\[
|A_{\Sigma}|\leq C\mathrm{e}^{N}/R.
\]

现固定 $x\in\Sigma^{0,2\pi}_{R,2R}$. 由于由引理~\ref{lem:1.36} 多值图是稳定的,
我们可以在稳定性不等式中使用 $\phi$, 得到
\[
\int_{\Sigma^{-\pi,3\pi}_{R/2,5R/2}}|A|^{2}
 \leq\int_\Sigma|A_{\Sigma}|^{2}\phi^{2}\leq\int|\nabla\phi|^{2}
 \leq\epsilon_E,
\]
其中最后一个不等号用到了 \eqref{eq:2.168}. 于是可以应用平均值不等式, 得到
\begin{equation}
\sup_{B_{3R/8}^{\Sigma}(x)}
 \bigl(R^{2}|\nabla A_{\Sigma}|^{2}+|A_{\Sigma}|^{2}\bigr)
 \leq C\epsilon_ER^{-2}.
\label{eq:2.169}
\end{equation}

由于 $\Sigma$ 与 $w$ 的图(局部地)都是有界的梯度图, 容易看出
\begin{equation}
\sup_{\Sigma^{-N\pi,N\pi}_{\mathrm{e}^{-N}R,\mathrm{e}^{N}R}}|\nabla w|
 \leq C\mathrm{e}^{N}\sup_\Sigma|w|/R
 \leq C\mathrm{e}^{N}\delta_p.
\label{eq:2.170}
\end{equation}
置 $v=\log w$. 选取 $\delta_p>0$(只依赖于 $N$), 使得 \eqref{eq:2.170} 蕴含 $w$ 在
$\Sigma^{-N\pi,N\pi}_{\mathrm{e}^{-N}R,\mathrm{e}^{N}R}$ 上满足 \eqref{eq:2.137}, 且
\[
|a|,\ |b|/|A|,\ |c|\leq1/4.
\]
对
\[
\dive\bigl(\phi^{2}\nabla v-\phi^{2}a\nabla v\bigr)
\]
应用 Stokes 定理并利用吸收不等式, 得到
\begin{equation}
\int_{B_{R/2}^{\Sigma}(x)}|\nabla v|^{2}
 \leq\int\phi^{2}|\nabla v|^{2}\leq CE(\phi)\leq C\epsilon_E.
\label{eq:2.171}
\end{equation}
结合 \eqref{eq:2.137} 与 \eqref{eq:2.169}, 一个简单的计算(如同 \eqref{eq:2.167} 中那样)表明: 在
$B_{3R/8}^{\Sigma}(x)$ 上
\begin{equation}
\Delta|\nabla v|^{2}
 \geq-C|\nabla v|^{4}-C\epsilon_ER^{-2}|\nabla v|^{2}
 -C\epsilon_ER^{-4}.
\label{eq:2.172}
\end{equation}
由证明定理~\ref{thm:2.2} 所用的伸缩论证, \eqref{eq:2.171} 与 \eqref{eq:2.172} 蕴含
$B_{R/4}^{\Sigma}(x)$ 上 $|\nabla v|^{2}$ 的逐点界; 把它与 \eqref{eq:2.169} 结合, 当
$\epsilon_E$ 很小时便得到 \eqref{eq:2.155}.

为得到 \eqref{eq:2.156}, 沿每条射线 $\theta=\theta_0$ 对
$|\nabla\log|w||\leq\alpha/\rho$ 积分, 得到
\begin{equation}
\log\frac{w(r_2,\theta_0)}{w(r_1,\theta_0)}
 \leq\alpha\int_{r_1}^{r_2}\rho^{-1}d\rho
 =\log\Bigl(\frac{r_2}{r_1}\Bigr)^{\alpha},
\label{eq:2.173}
\end{equation}
然后利用初等不等式
\[
(s-R)/R\leq2\log(s/R).
\]
\end{proof}

% ================================================================
\section{极小锥}
\label{sec:2.9}

我们用欧氏空间中极小锥的讨论来结束本章. 对这些锥的研究, 无论在 Bernstein
定理的推广中, 还是在局部正则性问题中, 都一直很重要.

若 $N^{k-1}$ 是 $\bS^{n-1}\subset\R^n$ 的(光滑)子流形, 则 $N$ 上的锥记作
$C(N)$, 作为集合
\begin{equation}
C(N)=\bigl\{x\in\R^n\mid x/|x|\in N\bigr\}.
\label{eq:2.174}
\end{equation}
因此 $C(N)$ 是去掉原点(即锥的顶点)后的光滑 $k$ 维子流形. 按定义, 锥在关于
原点的伸缩下不变. 首先我们有下面这个简单的引理, 其证明留给读者:

\begin{lemma}
\label{lem:2.36}
子流形 $N^{k-1}\subset\bS^{n-1}$ 是极小的, 当且仅当它的欧氏平均曲率处处与
$\bS^{n-1}\subset\R^n$ 正交.
\end{lemma}

设 $x=(x_1,\dots,x_n)$ 表示 $\R^n$ 中的位置向量, 并置 $r=|x|$. 由该引理可知,
若 $N^{k-1}\subset\bS^{n-1}$ 极小, 则 $N$ 的欧氏平均曲率由向量场
\begin{equation}
\Delta x=(\Delta x_1,\dots,\Delta x_n)
\label{eq:2.175}
\end{equation}
给出, 其中 $\Delta$ 是 $N$ 上的度量 Laplace 算子.

现在我们得到下述引理:

\begin{lemma}
\label{lem:2.37}
若 $N^{k-1}\subset\bS^{n-1}$ 是极小子流形, 则坐标函数是以 $k-1$ 为特征值的
特征函数.
\end{lemma}

\begin{proof}
由于 $N$ 极小, $\Delta x$ 与 $\bS^{n-1}$ 正交, 从而与位置向量 $x$ 同向.
于是存在函数 $f$, 使得
\begin{equation}
\Delta x=xf.
\label{eq:2.176}
\end{equation}
由于 $x\in\bS^{n-1}$, $|x|^{2}=1$, 故
\begin{equation}
0=\Delta|x|^{2}=2\langle x,\Delta x\rangle+2|\nabla x|^{2}
 =2f+2(k-1),
\label{eq:2.177}
\end{equation}
于是函数 $f$ 等于常数 $(1-k)$, 引理得证.
\end{proof}

\begin{example}
\label{ex:2.38}
若 $\bS^{k-1}$ 是 $\bS^{n-1}$ 中的全测地 $(k-1)$ 维球面, 则
$C(\bS^{k-1})$ 是 $\R^n$ 中过原点的 $k$ 维平面.
\end{example}

设 $N^{k-1}\subset\bS^{n-1}$ 是浸入子流形(不必极小). 由于
\begin{equation}
\nabla_{C(N)}u=r^{-1}\nabla_Nu(r^{-1}\,\cdot\,)
 +\frac{\partial u}{\partial r}\frac{\partial}{\partial r},
\label{eq:2.178}
\end{equation}
直接计算表明, $N$ 与 $C(N)$ 的 Laplace 在 $x\neq0$ 处满足下述简单公式:
\begin{equation}
\Delta_{C(N)}u=r^{-2}\Delta_Nu(r^{-1}x)
 +(k-1)r^{-1}\frac{\partial}{\partial r}u
 +\frac{\partial^{2}}{\partial r^{2}}u.
\label{eq:2.179}
\end{equation}

\begin{lemma}
\label{lem:2.39}
若 $N^{k-1}\subset\bS^{n-1}$ 是极小子流形, 则 $C(N)\subset\R^n$ 极小.
\end{lemma}

\begin{proof}
只需证明每个坐标函数 $x_i$ 在 $C(N)$ 上调和. 我们写成 $x_i=ru_i$, 其中 $u_i$
与 $r$ 无关, 且在 $N\subset\bS^{n-1}$ 上 $x_i$ 与 $u_i$ 一致. 引理~\ref{lem:2.37}
与 \eqref{eq:2.179} 给出
\begin{equation}
\begin{aligned}
\Delta_{C(N)}x_i
&=r^{-1}\Delta_Nu_i+u_i(k-1)r^{-1}\frac{\partial}{\partial r}r
 +u_i\frac{\partial^{2}}{\partial r^{2}}r\\
&=-(k-1)r^{-1}u_i+(k-1)r^{-1}u_i=0.
\end{aligned}
\label{eq:2.180}
\end{equation}
\end{proof}

极小锥出现在 Bernstein 原定理的各种推广中, 也出现在极小子流形的局部正则性
研究中. 这里的关键事实是: 由于锥在伸缩下不变, 比值
\begin{equation}
\frac{\Vol\bigl(C(N)\cap B_r\bigr)}{r^{k}}
\label{eq:2.181}
\end{equation}
是常数.

反之, 若 $\Sigma^{k}$ 极小且这个比值是常数, 则单调性公式中取到等号. 由此
$\nabla|x|^{2}$ 几乎处处与 $\Sigma$ 相切, 从而 $\Sigma$ 在伸缩下不变.

现在设 $\Sigma^{k}$ 是面积最小化极小超曲面. 像第 1 章中对极小图在 \eqref{eq:1.22} 处
所做的那样论证, 密度 $V_{\Sigma}$ 一致有界. 若 $r_j\to0$, 则伸缩后的曲面序列
$\Sigma_j=r_j\Sigma$ 也是面积最小化的. $V_{\Sigma_j}$ 的界蕴含存在收敛子列
(比如说, 作为面积最小化流或驻定变分流形收敛); 关于这一点下一章有更多讨论.
这个(未必唯一的)极限记作 $\Sigma_\infty$, 它也是面积最小化的. 由单调性公式
可知 $V_{\Sigma_\infty}=V_{\Sigma}$, 它也等于 $\Sigma_\infty$ 在原点的密度.
因此单调性公式中取到等号, 从而 $\Sigma_\infty$ 是面积最小化极小锥. 这一论证
可以追溯到 W. Fleming \cite{Fl}, 它使人能够通过证明面积最小化锥不存在来证明
Bernstein 定理.

\begin{theorem}[Fleming, \cite{Fl}]
\label{thm:2.40}
若 $\R^n$ 上极小曲面方程存在非线性的整解, 则 $\R^{n+1}$ 中存在奇异的面积
最小化锥.
\end{theorem}

Fleming 的结果发表后不久, E. De Giorgi 在 \cite{DG} 中证明: 定理~\ref{thm:2.40}
给出的极小锥实际上是柱状的. 这里``柱状''是指, 该锥作为 $\R^n$ 的子集等距地
分裂出一条直线, 即
\begin{equation}
C(N)=\R\times C(N')\subset\R^n.
\label{eq:2.182}
\end{equation}
这就蕴含 $\R^n$ 中确实存在奇异的面积最小化锥.

这一方法的可贵之处在于: 锥要容易分析得多. 这一联系是现代正则性理论的基石;
一个重要例证就是 Federer 所谓的维数下降论证. 例如, 在二维情形, 这样的锥
必定是过原点的一族平面; 容易看出, 只有当只有一个平面时它才可能是极小化的.
F. J. Almgren, Jr. 证明了在三维情形这样的锥也必定是超平面. 最后, J. Simons
对 $n\leq7$ 证明了同一定理.

\begin{theorem}[F. J. Almgren, Jr. \cite{Am1}(当 $n=3$), J. Simons \cite{Sim}
(当 $n\leq7$)]
\label{thm:2.41}
当 $n\leq7$ 时, 超平面是 $\R^n$ 中唯一的稳定极小超锥.
\end{theorem}

然而, 1969 年 Bombieri、De Giorgi 与 Giusti \cite{BDGG} 给出了 $\R^8$ 中
一个面积最小化奇异锥的例子. 事实上, 他们证明了: 当 $m\geq4$ 时, 锥
\begin{equation}
C_m=\bigl\{(x_1,\dots,x_{2m})\mid x_1^{2}+\cdots+x_m^{2}
 =x_{m+1}^{2}+\cdots+x_{2m}^{2}\bigr\}\subset\R^{2m}
\label{eq:2.183}
\end{equation}
是面积最小化的.
